% ============================================================
% HSBI --- Quantitative Research Methods
% FINAL Mock Examination --- SET B (Lectures 1--12, comprehensive)
% Focus: ISLP Chapters 7, 8, 10, 13 (~60%) + Chapters 2--6 (~40%)
% Prof. Dr. Christoph Weisser --- Summer Semester 2026
% ------------------------------------------------------------
% SINGLE SOURCE with solutions toggle:
%   Exam only:
%     pdflatex -jobname=Final_Mock_Exam_B final_mock_exam_b.tex
%   With solutions:
%     pdflatex -jobname=Final_Mock_Exam_B_Solutions "\def\withsolutions{1}% ============================================================
% HSBI --- Quantitative Research Methods
% FINAL Mock Examination --- SET B (Lectures 1--12, comprehensive)
% Focus: ISLP Chapters 7, 8, 10, 13 (~60%) + Chapters 2--6 (~40%)
% Prof. Dr. Christoph Weisser --- Summer Semester 2026
% ------------------------------------------------------------
% SINGLE SOURCE with solutions toggle:
%   Exam only:
%     pdflatex -jobname=Final_Mock_Exam_B final_mock_exam_b.tex
%   With solutions:
%     pdflatex -jobname=Final_Mock_Exam_B_Solutions "\def\withsolutions{1}% ============================================================
% HSBI --- Quantitative Research Methods
% FINAL Mock Examination --- SET B (Lectures 1--12, comprehensive)
% Focus: ISLP Chapters 7, 8, 10, 13 (~60%) + Chapters 2--6 (~40%)
% Prof. Dr. Christoph Weisser --- Summer Semester 2026
% ------------------------------------------------------------
% SINGLE SOURCE with solutions toggle:
%   Exam only:
%     pdflatex -jobname=Final_Mock_Exam_B final_mock_exam_b.tex
%   With solutions:
%     pdflatex -jobname=Final_Mock_Exam_B_Solutions "\def\withsolutions{1}% ============================================================
% HSBI --- Quantitative Research Methods
% FINAL Mock Examination --- SET B (Lectures 1--12, comprehensive)
% Focus: ISLP Chapters 7, 8, 10, 13 (~60%) + Chapters 2--6 (~40%)
% Prof. Dr. Christoph Weisser --- Summer Semester 2026
% ------------------------------------------------------------
% SINGLE SOURCE with solutions toggle:
%   Exam only:
%     pdflatex -jobname=Final_Mock_Exam_B final_mock_exam_b.tex
%   With solutions:
%     pdflatex -jobname=Final_Mock_Exam_B_Solutions "\def\withsolutions{1}\input{final_mock_exam_b.tex}"
% ============================================================
\documentclass[11pt,a4paper]{article}
\usepackage[utf8]{inputenc}
\usepackage[T1]{fontenc}
\usepackage[english]{babel}
\usepackage[margin=2.2cm]{geometry}
\usepackage{amsmath,amssymb}
\usepackage{booktabs,array}
\usepackage{enumitem}
\usepackage{xcolor}
\usepackage{tcolorbox}
\tcbuselibrary{skins,breakable}

\setlength{\parindent}{0pt}

% Teal solution box (only shown when \withsolutions is defined)
\newtcolorbox{solutionbox}{enhanced,breakable,colback=teal!5,colframe=teal!55!black,
  boxrule=0pt,leftrule=2.5pt,arc=2pt,fonttitle=\bfseries,title=Solution,
  left=2mm,right=2mm,top=1mm,bottom=1.5mm}

\newcommand{\problem}[2]{\par\vspace{7mm}\noindent{\large\textbf{Problem #1}}%
  \hfill{\large\textbf{(#2 points)}}\par\nopagebreak\vspace{0.6mm}%
  \noindent\rule{\textwidth}{0.6pt}\par\nopagebreak\vspace{2.5mm}}
\newcommand{\pp}[1]{\textbf{[#1~P]}\enspace}

\begin{document}

% ------------------------------------------------------------
% Header block
% ------------------------------------------------------------
\begin{center}
  {\Large\textbf{HSBI --- Bielefeld University of Applied Sciences and Arts}}\\[1.2mm]
  {\large Quantitative Research Methods}\\[1mm]
  {\normalsize Summer Semester 2026 \quad$\cdot$\quad Examiner: Prof.\ Dr.\ Christoph Weisser}\\[4.5mm]
  {\LARGE\textbf{Final Mock Examination --- Set B (Lectures 1--12)}}\\[1.2mm]
  {\normalsize Comprehensive coverage of ISLP Chapters 2--8, 10 and 13,\\
   with emphasis on Chapters 7 (splines/GAMs), 8 (trees), 10 (deep learning) and 13 (multiple testing)}%
  \ifdefined\withsolutions\\[2.2mm]{\large\textcolor{teal!55!black}{\textbf{--- Version with detailed solutions ---}}}\fi
\end{center}

\vspace{2.5mm}
\begin{center}
\begin{tabular}{@{}ll@{}}
  \textbf{Duration:} & \textbf{120 minutes}\\
  \textbf{Total credit:} & \textbf{120 points}\\
  \textbf{Allowed aids:} & non-programmable calculator; one handwritten A4 sheet of notes (both sides)\\
\end{tabular}
\end{center}

\vspace{3mm}
Name: \rule[-1mm]{0.38\textwidth}{0.4pt}\hfill Matriculation no.: \rule[-1mm]{0.24\textwidth}{0.4pt}

\vspace{4.5mm}
\begin{center}
\renewcommand{\arraystretch}{1.35}
\begin{tabular}{|l|c|c|c|c|c|c|c|c|}
\hline
\textbf{Problem} & \textbf{P1} & \textbf{P2} & \textbf{P3} & \textbf{P4} & \textbf{P5} & \textbf{P6} & \textbf{P7} & \textbf{Total}\\\hline
Maximum points & 16 & 8 & 20 & 16 & 20 & 20 & 20 & \textbf{120}\\\hline
Achieved points & \hspace{8mm} & \hspace{8mm} & \hspace{8mm} & \hspace{8mm} & \hspace{8mm} & \hspace{8mm} & \hspace{8mm} & \hspace{10mm}\\\hline
\end{tabular}
\end{center}

\vspace{2.5mm}
\textbf{Instructions}
\begin{itemize}[nosep,leftmargin=5.5mm]
  \item Justify all answers. Numerical results without a derivation or a brief reasoning receive no credit.
  \item Round final numerical results to 3 decimal places unless stated otherwise. Small deviations caused by carrying rounded intermediate results are accepted.
  \item All numerical constants you need (powers of $e$, logarithms, $(0.95)^{10}$, $t$-quantiles) are provided in the problems.
  \item The number of points per problem roughly equals the recommended working time in minutes.
\end{itemize}
\vspace{1mm}\noindent\rule{\textwidth}{0.6pt}

% ============================================================
\problem{1: Moving beyond linearity --- splines}{16}
{\small\itshape Lecture slides: Chapter 7 --- Moving Beyond Linearity (Lecture 9).}\par\vspace{1mm}

Throughout this problem, consider regression splines built on a single predictor $X$ using
$K=5$ interior knots.

\begin{enumerate}[label=(\alph*),leftmargin=7mm,itemsep=2.5mm]
  \item \pp{4} A \emph{cubic spline} is a piecewise cubic polynomial that is continuous and has
    continuous first and second derivatives at every knot. How many degrees of freedom
    (free parameters, counting the intercept) does a cubic spline with $K$ interior knots use?
    Justify your answer by counting the parameters of the polynomial pieces and subtracting the
    constraints at the knots, and evaluate for $K=5$.
  \item \pp{4} How many degrees of freedom does a \emph{natural} cubic spline with the same
    $K=5$ knots use? State the additional constraint a natural spline imposes, in which region
    of the predictor space it acts, and why practitioners find this constraint useful.
  \item \pp{4} A \emph{smoothing spline} is the function $g$ that minimises
    \[
      \sum_{i=1}^{n}\bigl(y_i-g(x_i)\bigr)^2+\lambda\int g''(t)^2\,dt .
    \]
    Explain the meaning of the two terms, and describe the shape of the solution $\hat{g}$ in
    the two limiting cases $\lambda\to 0$ and $\lambda\to\infty$, giving the effective degrees
    of freedom that $\hat{g}$ approaches in each case.
  \item \pp{4} Write down the \emph{truncated-power basis} representation of a cubic spline with
    a single knot at $\xi$, define the truncated-power basis function $(x-\xi)^3_+$ explicitly,
    and explain in one sentence why adding this term keeps the fit continuous with continuous
    first and second derivatives at $\xi$.
\end{enumerate}

\ifdefined\withsolutions
\begin{solutionbox}
\textbf{(a) Degrees of freedom of a cubic spline.}
With $K$ interior knots the range of $X$ is divided into $K+1$ intervals, and on each interval
the fit is a cubic with $4$ coefficients, giving $4(K+1)=4K+4$ free parameters. At every knot
the requirement that $g$, $g'$ and $g''$ join continuously imposes $3$ constraints, i.e.\ $3K$
constraints in total. Hence
\[
  \mathrm{df}=4(K+1)-3K=K+4 .
\]
For $K=5$: $\mathrm{df}=5+4=\mathbf{9}$ (equivalently the basis intercept, $x$, $x^2$, $x^3$
plus one truncated-power term per knot, $4+K$).

\medskip
\textbf{(b) Natural cubic spline.}
A natural spline adds the requirement that the function be \emph{linear beyond the two boundary
knots}, i.e.\ on the outer regions $X<\xi_1$ and $X>\xi_K$. Forcing the quadratic and cubic
terms to vanish in each of those two regions removes $2$ parameters per side, $4$ in total:
\[
  \mathrm{df}=(K+4)-4=K,\qquad K=5:\ \mathrm{df}=\mathbf{5}.
\]
Usefulness: ordinary polynomials and unconstrained splines behave erratically near the
boundaries, where the data are sparse; forcing linearity there reduces the \emph{variance} of
predictions at and beyond the outer knots, at the cost of only a little extra bias.

\medskip
\textbf{(c) Smoothing spline.}
The first term (the RSS) rewards closeness to the observed data, while the second term
penalises \emph{roughness}: $g''$ measures the local curvature of $g$, so $\int g''(t)^2dt$ is
large for wiggly functions, and $\lambda\ge 0$ controls the trade-off.
\emph{$\lambda\to 0$:} the penalty vanishes, so $\hat{g}$ may be arbitrarily wiggly and
\emph{interpolates} every training point (zero training RSS, extreme overfitting); the
effective degrees of freedom approach $n$.
\emph{$\lambda\to\infty$:} any curvature becomes infinitely costly, forcing $g''\equiv 0$, so
$\hat{g}$ is a straight line, namely the ordinary \emph{least-squares regression line}
(effective df $\to 2$). For intermediate $\lambda$ (chosen e.g.\ by LOOCV) the effective df
decrease smoothly from $n$ down to $2$.

\medskip
\textbf{(d) Truncated-power basis, single knot $\xi$.}
\[
  f(x)=\beta_0+\beta_1x+\beta_2x^2+\beta_3x^3+\beta_4\,(x-\xi)^3_+,
  \qquad
  (x-\xi)^3_+=\begin{cases}(x-\xi)^3, & x>\xi\\[0.5mm] 0, & x\le\xi\end{cases}
\]
At $x=\xi$ the added function $(x-\xi)^3_+$ and its first and second derivatives are all $0$, so
including it alters only the \emph{third} derivative across the knot; the fit therefore stays
continuous with continuous first and second derivatives --- precisely the smoothness that
defines a cubic spline.
\end{solutionbox}
\fi

% ============================================================
\problem{2: Generalised additive models}{8}
{\small\itshape Lecture slides: Chapter 7 --- Moving Beyond Linearity (Lecture 9).}\par\vspace{1mm}

\begin{enumerate}[label=(\alph*),leftmargin=7mm,itemsep=2.5mm]
  \item \pp{3} Write down the general form of a \emph{generalised additive model} (GAM) for a
    quantitative response $Y$ and predictors $X_1,\dots,X_p$, and explain in one sentence how it
    generalises multiple linear regression.
  \item \pp{3} State one important \emph{interpretability advantage} of GAMs (concerning the
    fitted functions $f_j$) and their central structural \emph{limitation}, including how that
    limitation can be repaired.
  \item \pp{2} Describe in one or two sentences how the \emph{backfitting} algorithm fits a GAM.
\end{enumerate}

\ifdefined\withsolutions
\begin{solutionbox}
\textbf{(a) Model form.}
\[
  Y=\beta_0+f_1(X_1)+f_2(X_2)+\dots+f_p(X_p)+\varepsilon .
\]
Multiple linear regression is the special case in which every $f_j(X_j)=\beta_jX_j$; a GAM
replaces each straight-line term by a flexible smooth function $f_j$ (e.g.\ a natural or
smoothing spline, or a local regression) while retaining the \emph{additive} structure.

\medskip
\textbf{(b) Advantage and limitation.}
\emph{Advantage (interpretability):} because the effects add up, the contribution of each
predictor is carried by its own function $f_j$, which can be plotted and read off one panel at
a time with the other predictors held fixed --- non-linear effects stay as easy to inspect as
coefficients in a linear model.
\emph{Limitation:} pure additivity rules out \emph{interactions} --- the effect of $X_j$ is not
allowed to depend on the value of another $X_k$. Repair: introduce interaction terms by hand,
e.g.\ a product $X_j\times X_k$ or a low-dimensional smooth surface $f_{jk}(X_j,X_k)$ as an
extra additive component.

\medskip
\textbf{(c) Backfitting.}
Loop over the predictors and repeatedly re-estimate one $f_j$ at a time (with a smoother) using
the \emph{partial residuals} $r_i=y_i-\beta_0-\sum_{k\ne j}f_k(x_{ik})$, keeping the other
fitted functions fixed; repeat the cycles until the $f_j$ no longer change appreciably.
\end{solutionbox}
\fi

% ============================================================
\problem{3: Regression and classification trees by hand}{20}
{\small\itshape Lecture slides: Chapter 8 --- Tree-Based Methods (Lecture 10).}\par\vspace{1mm}

A logistics firm records, for $n=6$ regional depots, the number of delivery vans $x$ and the
number of parcels delivered per day $y$ (in thousands):

\begin{center}
\begin{tabular}{lcccccc}
\toprule
Depot $i$ & 1 & 2 & 3 & 4 & 5 & 6\\
\midrule
Vans $x_i$   & 10 & 12 & 14 & 16 & 18 & 20\\
Parcels $y_i$ & 6  & 4  & 5  & 11 & 13 & 12\\
\bottomrule
\end{tabular}
\end{center}

A regression tree with a \emph{single} split at cutpoint $s$ partitions the depots into
$R_1=\{i: x_i<s\}$ and $R_2=\{i: x_i\ge s\}$ and predicts the region mean $\hat{y}_{R_m}$ in
each region.

\begin{enumerate}[label=(\alph*),leftmargin=7mm,itemsep=2.5mm]
  \item \pp{8} For each of the two candidate cutpoints $s=15$ and $s=17$, determine the two
    regions, their means, and the total residual sum of squares
    $\operatorname{RSS}(s)=\sum_{m=1}^{2}\sum_{i\in R_m}(y_i-\hat{y}_{R_m})^2$.
    Which cutpoint does recursive binary splitting prefer, and why?
  \item \pp{2} Suppose recursive binary splitting selects the split at $s=15$, giving regions
    $R_1=\{x<15\}$ with mean $\hat{y}_{R_1}=5$ and $R_2=\{x\ge 15\}$ with mean
    $\hat{y}_{R_2}=12$. Predict the number of parcels for a new depot with $x=13$ vans.
  \item \pp{6} A \emph{classification} tree for a different task has a terminal node containing
    $12$ training observations from three classes: $6$ ``Premium'', $3$ ``Standard'' and $3$
    ``Basic''. Compute the class proportions, the Gini index
    $G=\sum_{k}\hat{p}_{k}(1-\hat{p}_{k})$ and the entropy $D=-\sum_{k}\hat{p}_{k}\ln\hat{p}_{k}$
    of this node \emph{(use $\ln 0.5=-0.6931$ and $\ln 0.25=-1.3863$)}. Which class does the node
    predict?
  \item \pp{4} In \emph{cost-complexity pruning} one minimises, over subtrees $T\subseteq T_0$
    of the large tree $T_0$,
    \[
      \sum_{m=1}^{|T|}\ \sum_{i:\,x_i\in R_m}\bigl(y_i-\hat{y}_{R_m}\bigr)^2+\alpha\,|T| .
    \]
    Explain the role of $\alpha$: what happens for $\alpha=0$, what happens as $\alpha$ grows,
    and how is $\alpha$ chosen in practice?
\end{enumerate}

\ifdefined\withsolutions
\begin{solutionbox}
\textbf{(a) Evaluating the two cutpoints.}
For reference the root mean is $\bar{y}=(6+4+5+11+13+12)/6=51/6=8.5$.

\emph{Cutpoint $s=15$:} $R_1=\{1,2,3\}$ with $y$-values $\{6,4,5\}$, mean $\hat{y}_{R_1}=15/3=5$;
$R_2=\{4,5,6\}$ with $y$-values $\{11,13,12\}$, mean $\hat{y}_{R_2}=36/3=12$.
\[
  \begin{aligned}
  \operatorname{RSS}_{R_1}&=(6-5)^2+(4-5)^2+(5-5)^2=1+1+0=2,\\
  \operatorname{RSS}_{R_2}&=(11-12)^2+(13-12)^2+(12-12)^2=1+1+0=2,
  \end{aligned}
\]
so $\operatorname{RSS}(15)=2+2=\mathbf{4}$.

\emph{Cutpoint $s=17$:} $R_1=\{1,2,3,4\}$ with $y$-values $\{6,4,5,11\}$, mean
$\hat{y}_{R_1}=26/4=6.5$; $R_2=\{5,6\}$ with $y$-values $\{13,12\}$, mean $\hat{y}_{R_2}=12.5$.
\[
  \operatorname{RSS}_{R_1}=(6-6.5)^2+(4-6.5)^2+(5-6.5)^2+(11-6.5)^2=0.25+6.25+2.25+20.25=29,
\]
\[
  \operatorname{RSS}_{R_2}=(13-12.5)^2+(12-12.5)^2=0.25+0.25=0.5,
\]
so $\operatorname{RSS}(17)=29+0.5=\mathbf{29.5}$.

Recursive binary splitting is \emph{greedy} and takes the cutpoint with the smaller RSS: since
$4<29.5$, it prefers $\mathbf{s=15}$, which cleanly separates the three low-throughput depots
(mean $5$) from the three high-throughput depots (mean $12$).

\medskip
\textbf{(b) Prediction for $x=13$.}
$x=13<15$, so the new depot falls into $R_1$ and the tree predicts
$\hat{y}=\hat{y}_{R_1}=\mathbf{5}$ (i.e.\ $5{,}000$ parcels per day).

\medskip
\textbf{(c) Node impurity (three classes).}
$\hat{p}_{\text{Prem}}=6/12=0.5$, $\hat{p}_{\text{Std}}=3/12=0.25$,
$\hat{p}_{\text{Basic}}=3/12=0.25$.
\[
  G=0.5(1-0.5)+0.25(1-0.25)+0.25(1-0.25)=0.25+0.1875+0.1875=\mathbf{0.625}.
\]
\[
  \begin{aligned}
  D&=-\bigl(0.5\ln 0.5+0.25\ln 0.25+0.25\ln 0.25\bigr)\\
   &=-\bigl(0.5(-0.6931)+0.5(-1.3863)\bigr)=0.3466+0.6931=\mathbf{1.040}.
  \end{aligned}
\]
The node predicts the majority class, \textbf{Premium} (estimated probability $0.5$). Both $G$
and $D$ would be $0$ for a pure node; here the node is quite impure, consistent with the split
carrying three classes.

\medskip
\textbf{(d) Role of $\alpha$.}
$\alpha\ge 0$ is the price charged per terminal node, trading the subtree's fit (first term)
against its size $|T|$.
For $\alpha=0$ the criterion reduces to the training RSS, so the full unpruned tree $T_0$ is
optimal. As $\alpha$ increases, each extra leaf must ``earn its keep''; a branch whose RSS
reduction is smaller than $\alpha$ is pruned away, so the optimal subtree shrinks and one
obtains a \emph{nested sequence} of subtrees as $\alpha$ climbs (weakest-link pruning). In
practice $\alpha$ is selected by \emph{cross-validation} --- pick the $\alpha$ giving the lowest
CV error, then refit that subtree on the full data --- which balances bias and variance rather
than merely minimising training error.
\end{solutionbox}
\fi

% ============================================================
\problem{4: Ensembles --- bagging, random forests, boosting}{16}
{\small\itshape Lecture slides: Chapter 8 --- Tree-Based Methods (Lecture 10).}\par\vspace{1mm}

\begin{enumerate}[label=(\alph*),leftmargin=7mm,itemsep=2.5mm]
  \item \pp{6} Let $\hat{f}_1(x),\dots,\hat{f}_B(x)$ be $B$ identically distributed (but not
    independent) tree predictions at a fixed point $x$, each with variance $\sigma^2$ and
    pairwise correlation $\rho$. The variance of the bagged average is
    \[
      \operatorname{Var}\Bigl(\tfrac{1}{B}\sum_{b=1}^{B}\hat{f}_b(x)\Bigr)
      =\rho\,\sigma^{2}+\frac{1-\rho}{B}\,\sigma^{2}.
    \]
    (i) Explain briefly why this formula shows that bagging reduces variance and which term
    limits the reduction. (ii) Evaluate the variance for $B=500$, $\rho=0.7$ and $\sigma^2=4$,
    and compare it with the variance of a single tree. (iii) What value does the variance
    approach as $B\to\infty$?
  \item \pp{4} How do \emph{random forests} modify bagging at each split, and what is the
    recommended subset size $m$ for \emph{classification} as a function of $p$ (evaluate for
    $p=36$)? Using the variance formula
    $\operatorname{Var}=\rho\sigma^2+\tfrac{1-\rho}{B}\sigma^2$ from part~(a), explain why this
    modification improves on bagging.
  \item \pp{4} A data scientist fits a boosting model for regression trees in
    \texttt{scikit-learn}:
    \begin{center}
    \texttt{GradientBoostingRegressor(n\_estimators=500, learning\_rate=0.01, max\_depth=1)}
    \end{center}
    These three arguments are exactly the tuning parameters of boosting. State in one sentence
    each what \texttt{n\_estimators}, \texttt{learning\_rate} and \texttt{max\_depth} control,
    and explain the ``slow learning'' idea behind the small \texttt{learning\_rate}.
  \item \pp{2} What does the \emph{out-of-bag} (OOB) error of a bagged model estimate, and why
    is it essentially free to obtain?
\end{enumerate}

\ifdefined\withsolutions
\begin{solutionbox}
\textbf{(a) Variance of the bagged average.}
(i) The second term $\frac{1-\rho}{B}\sigma^2$ shrinks to $0$ as $B$ grows, so averaging over
many trees removes the ``independent part'' of their variability --- this is the variance
reduction of bagging. The first term $\rho\sigma^2$ does \emph{not} depend on $B$: the
bootstrap trees are grown on overlapping resamples of the same data (and share the same strong
predictors), so they are positively correlated, and this common part of the variance cannot be
averaged away. It is the floor that limits the reduction.
(ii) For $B=500$, $\rho=0.7$, $\sigma^2=4$:
\[
  \operatorname{Var}=0.7\cdot 4+\frac{1-0.7}{500}\cdot 4
  =2.8+\frac{0.3\cdot 4}{500}=2.8+0.0024=\mathbf{2.802},
\]
versus $\sigma^2=4$ for a single tree --- a reduction of about $30\%$.
(iii) As $B\to\infty$ the variance approaches the floor $\rho\sigma^2=0.7\cdot 4=\mathbf{2.8}$.
Enlarging $B$ further does not overfit; it only stabilises the average.

\medskip
\textbf{(b) Random forests.}
At each split the tree may consider only a fresh \emph{random subset} of $m$ out of the $p$
predictors as candidates. For classification the usual choice is $m\approx\sqrt{p}$, so for
$p=36$: $m=\sqrt{36}=\mathbf{6}$.
Why it helps: in plain bagging a dominant predictor is chosen at the top split of nearly every
tree, making the trees look alike and pushing $\rho$ up. Withholding that predictor from about
$(p-m)/p$ of the split searches \emph{decorrelates} the trees, i.e.\ lowers $\rho$; by (a) this
lowers the irreducible floor $\rho\sigma^2$ of the ensemble variance, which plain averaging can
never reach.

\medskip
\textbf{(c) Boosting tuning parameters (read from the call).}
\begin{itemize}[nosep,leftmargin=5mm]
  \item \texttt{n\_estimators}$=B=500$: the number of small trees added in sequence; unlike
    bagging, boosting \emph{can} overfit if $B$ is too large, so $B$ is tuned by cross-validation.
  \item \texttt{learning\_rate}$=\lambda=0.01$ (shrinkage): scales the contribution of each new
    tree and hence the speed of learning; a smaller $\lambda$ usually requires a larger $B$.
  \item \texttt{max\_depth}$=d=1$ (interaction depth): the number of splits per tree and the
    interaction order captured; $d=1$ means \emph{stumps} (an additive model).
\end{itemize}
\emph{Slow learning:} each new tree is fit to the current \emph{residuals}, and only the small
fraction \texttt{learning\_rate}$=\lambda$ of it is added to the model, so the fit improves
gradually and precisely where it is still weak, rather than fitting the data hard in a single
step --- which tends to generalise better.

\medskip
\textbf{(d) OOB error.}
Each bootstrap sample omits, on average, about one third of the observations
(``out-of-bag''). Predicting each observation only from the trees that did \emph{not} train on
it, and aggregating, gives the OOB error --- a valid estimate of the bagged model's
\emph{test} error. It is essentially free because it reuses a by-product of the bootstrap; no
separate validation set or cross-validation loop is needed.
\end{solutionbox}
\fi

% ============================================================
\problem{5: Neural networks by hand}{20}
{\small\itshape Lecture slides: Chapter 10 --- Deep Learning (Lecture 11).}\par\vspace{1mm}

Consider a small feed-forward network for a quantitative response with input
$x=(x_1,x_2)=(2,1)$, one hidden layer of two ReLU units ($g(z)=\max\{0,z\}$) and a linear
output unit:
\[
  A_1=g\bigl(-2+3\,x_1-1\,x_2\bigr),\qquad
  A_2=g\bigl(1-2\,x_1+1\,x_2\bigr),\qquad
  f(x)=-1+2\,A_1+3\,A_2 .
\]

\begin{enumerate}[label=(\alph*),leftmargin=7mm,itemsep=2.5mm]
  \item \pp{8} Carry out the forward pass step by step: compute both pre-activations, both
    activations $A_1,A_2$ and the output $f(x)$. Comment briefly on what the ReLU did to the
    second unit and why non-linear activations are essential.
  \item \pp{6} A classification network with $K=3$ classes returns the logits (output-layer
    pre-activations) $z=(z_1,z_2,z_3)=(2.0,\ 1.0,\ -1.0)$ for one observation whose true class
    is class~1. Compute all three softmax probabilities and the cross-entropy loss
    $-\ln\hat{p}_{\text{true}}$.
    \emph{(Use $e^{2}=7.389$, $e^{1}=2.718$, $e^{-1}=0.368$ and $\ln 10.475=2.349$.)}
  \item \pp{4} Count the parameters of a fully connected network with $p=6$ inputs, one hidden
    layer of $k=4$ units and an output layer of $K=3$ units, where every hidden and output unit
    carries a bias. Show the count separately for the two weight layers.
  \item \pp{2} A convolutional layer receives a $32\times 32$ image and applies filters of size
    $F=5$ with padding $P=0$ and stride $S=1$. Compute the spatial size of the output feature
    map from $\bigl(W-F+2P\bigr)/S+1$.
\end{enumerate}

\ifdefined\withsolutions
\begin{solutionbox}
\textbf{(a) Forward pass.}
Pre-activations:
\[
  z_1=-2+3\cdot 2-1\cdot 1=-2+6-1=3,\qquad
  z_2=1-2\cdot 2+1\cdot 1=1-4+1=-2 .
\]
ReLU activations:
\[
  A_1=g(3)=\max\{0,3\}=3,\qquad A_2=g(-2)=\max\{0,-2\}=0 .
\]
Output:
\[
  f(x)=-1+2\cdot 3+3\cdot 0=-1+6+0=\mathbf{5}.
\]
The ReLU \emph{switched the second unit off}: its negative pre-activation is clipped to $0$, so
unit~2 contributes nothing for this particular input (though it could for others). This
thresholding is exactly the source of non-linearity: with a purely linear $g$, a network of any
depth would collapse to a single linear function of $x$.

\medskip
\textbf{(b) Softmax and cross-entropy.}
Exponentials $e^{2.0}=7.389$, $e^{1.0}=2.718$, $e^{-1.0}=0.368$ sum to
$7.389+2.718+0.368=10.475$. Softmax probabilities:
\[
  \hat{p}_1=\frac{7.389}{10.475}=0.705,\qquad
  \hat{p}_2=\frac{2.718}{10.475}=0.259,\qquad
  \hat{p}_3=\frac{0.368}{10.475}=0.035,
\]
which sum to $0.999\approx 1$. Cross-entropy loss for the true class~1:
\[
  -\ln\hat{p}_1=-\ln\frac{e^{2}}{10.475}=-\bigl(2-\ln 10.475\bigr)=2.349-2=\mathbf{0.349}.
\]
(Equivalently $-\ln 0.705\approx 0.349$; a perfect prediction $\hat{p}_1=1$ would give loss 0.)

\medskip
\textbf{(c) Parameter count for $6\to 4\to 3$.}
Input $\to$ hidden: $6\cdot 4=24$ weights $+\,4$ biases $=28$.
Hidden $\to$ output: $4\cdot 3=12$ weights $+\,3$ biases $=15$.
Total: $28+15=\mathbf{43}$ parameters.

\medskip
\textbf{(d) Convolution output size.}
\[
  \frac{W-F+2P}{S}+1=\frac{32-5+2\cdot 0}{1}+1=\frac{27}{1}+1=28,
\]
so the output feature map is $\mathbf{28\times 28}$. With no padding a $5\times 5$ filter at
stride $1$ shrinks each spatial dimension by $F-1=4$.
\end{solutionbox}
\fi

% ============================================================
\problem{6: Multiple testing}{20}
{\small\itshape Lecture slides: Chapter 13 --- Multiple Testing (Lecture 12).}\par\vspace{1mm}

\begin{enumerate}[label=(\alph*),leftmargin=7mm,itemsep=2.5mm]
  \item \pp{4} A lab runs $m=10$ \emph{independent} hypothesis tests, each at level
    $\alpha=0.05$, and suppose all $10$ null hypotheses are true. Define the family-wise error
    rate (FWER) and compute it here \emph{(use $(0.95)^{10}=0.5987$)}. Interpret the result in
    one sentence.
  \item \pp{7} A team screens $m=6$ candidate biomarkers; the ordered $p$-values are
    \begin{center}
    \begin{tabular}{lcccccc}
    \toprule
    $j$ & 1 & 2 & 3 & 4 & 5 & 6\\
    \midrule
    $p_{(j)}$ & 0.002 & 0.009 & 0.012 & 0.015 & 0.040 & 0.300\\
    \bottomrule
    \end{tabular}
    \end{center}
    Control the FWER at $\alpha=0.05$ (i) with \emph{Bonferroni} and (ii) with the
    \emph{Holm step-down} procedure. Give the relevant thresholds, state exactly which
    hypotheses each method rejects, and explain why Holm can never reject fewer hypotheses than
    Bonferroni.
  \item \pp{6} Apply the \emph{Benjamini--Hochberg} (BH) procedure at FDR level $q=0.05$ to the
    same six ordered $p$-values ($m=6$; $p_{(1)},\dots,p_{(6)}=0.002,\ 0.009,\ 0.012,\ 0.015,\
    0.040,\ 0.300$): give the thresholds $jq/m$, determine which hypotheses are rejected, and
    state precisely what ``FDR control at $q=0.05$'' guarantees for the resulting set of
    rejections.
  \item \pp{3} A pharmaceutical company runs a \emph{confirmatory} Phase-III trial with $m=4$
    pre-specified primary endpoints; a single significant endpoint could support regulatory
    approval of the drug. Should the company control the FWER or the FDR? Justify your
    recommendation.
\end{enumerate}

\ifdefined\withsolutions
\begin{solutionbox}
\textbf{(a) FWER for $m=10$ independent tests.}
$\text{FWER}=\Pr(\text{at least one Type I error})=\Pr(V\ge 1)$, where $V$ counts the true
nulls that are falsely rejected. With independence and all nulls true,
\[
  \text{FWER}=1-\Pr(\text{no false rejection})=1-(1-0.05)^{10}=1-(0.95)^{10}
  =1-0.5987=\mathbf{0.4013}.
\]
Interpretation: although each individual test uses the $5\%$ level, there is about a $40\%$
chance of at least one false discovery across the ten tests --- testing many hypotheses without
adjustment makes spurious findings very likely.

\medskip
\textbf{(b) Bonferroni and Holm at $\alpha=0.05$.}
\emph{Bonferroni:} reject $H_{(j)}$ iff $p_{(j)}\le\alpha/m=0.05/6=0.00833$.
Only $p_{(1)}=0.002\le 0.00833$ \checkmark; already $p_{(2)}=0.009>0.00833$ $\times$.
$\Rightarrow$ Bonferroni rejects $H_{(1)}$ alone: \textbf{1 rejection}. (Note $p_{(2)}=0.009$ is
below the naive $0.05$ cutoff, yet Bonferroni does not reject it.)

\emph{Holm:} compare $p_{(j)}$ with $\alpha/(m-j+1)$ in order and stop at the first failure:
\begin{center}
\begin{tabular}{lcccccc}
\toprule
$j$ & 1 & 2 & 3 & 4 & 5 & 6\\
\midrule
$p_{(j)}$ & 0.002 & 0.009 & 0.012 & 0.015 & 0.040 & 0.300\\
$\alpha/(m-j+1)$ & 0.00833 & 0.0100 & 0.0125 & 0.0167 & 0.0250 & 0.0500\\
reject? & \checkmark & \checkmark & \checkmark & \checkmark & stop & ---\\
\bottomrule
\end{tabular}
\end{center}
$0.002\le 0.00833$, $0.009\le 0.0100$, $0.012\le 0.0125$, $0.015\le 0.0167$, but
$0.040>0.0250$: stop. $\Rightarrow$ Holm rejects $H_{(1)},H_{(2)},H_{(3)},H_{(4)}$:
\textbf{4 rejections}.

\emph{Why Holm $\supseteq$ Bonferroni:} Holm's thresholds $\alpha/(m-j+1)$ are all $\ge\alpha/m$
(with equality only at $j=1$), so any $p$-value that clears the strict Bonferroni bar also
clears Holm's more lenient later hurdles --- Holm is uniformly at least as powerful, while still
controlling the FWER at $\alpha$ with no independence assumption.

\medskip
\textbf{(c) Benjamini--Hochberg at $q=0.05$.}
Thresholds $jq/m=j\cdot 0.05/6$:
\begin{center}
\begin{tabular}{lcccccc}
\toprule
$j$ & 1 & 2 & 3 & 4 & 5 & 6\\
\midrule
$p_{(j)}$ & 0.002 & 0.009 & 0.012 & 0.015 & 0.040 & 0.300\\
$jq/m$ & 0.00833 & 0.0167 & 0.0250 & 0.0333 & 0.0417 & 0.0500\\
$p_{(j)}\le jq/m$? & yes & yes & yes & yes & \textbf{yes} & no\\
\bottomrule
\end{tabular}
\end{center}
$L=$ largest $j$ with $p_{(j)}\le jq/m$: here $j=5$ (since $0.040\le 0.0417$, while
$0.300>0.0500$). BH rejects \emph{all} $H_{(j)}$ with $j\le L$, i.e.\
$H_{(1)},\dots,H_{(5)}$: \textbf{5 rejections}.
Guarantee: $\text{FDR}=E\bigl[V/\max(R,1)\bigr]\le q=0.05$ --- \emph{on average} at most $5\%$
of the rejected hypotheses are false discoveries. It does \emph{not} promise that at most $5\%$
of these particular $5$ rejections are false, nor does it bound the probability of making any
false rejection (that would be FWER control).

\medskip
\textbf{(d) Confirmatory trial recommendation.}
Control the \textbf{FWER}. Here even a \emph{single} false positive is very costly --- it could
lead to approving an ineffective drug --- so the study must guard against any Type I error, not
merely keep their average proportion small. With only $m=4$ pre-specified endpoints, an FWER
procedure (e.g.\ Holm or Bonferroni at $0.05/4=0.0125$) sacrifices little power. FDR control,
appropriate for large exploratory screens where a few false leads are tolerable, is too
permissive for a confirmatory, decision-driving trial.
\end{solutionbox}
\fi

% ============================================================
\problem{7: Cumulative essentials (Chapters 2--6)}{20}
{\small\itshape Lecture slides: Chapters 2--6 (Lectures 1--8).}\par\vspace{1mm}

\begin{enumerate}[label=(\alph*),leftmargin=7mm,itemsep=2.5mm]
  \item \pp{4} A medical screening classifier is evaluated on $200$ patients (label
    \texttt{1}$=$diseased, \texttt{0}$=$healthy). In \texttt{scikit-learn},
    \texttt{confusion\_matrix(y\_true, y\_pred)} returns the array below (rows $=$ true class,
    columns $=$ predicted class; label order $[0,1]$):
    \begin{center}
    \begin{tabular}{lcc}
    \toprule
     & predicted \texttt{0} & predicted \texttt{1}\\
    \midrule
    true \texttt{0} (healthy)  & $120$ & $30$\\
    true \texttt{1} (diseased) & $10$  & $40$\\
    \bottomrule
    \end{tabular}
    \end{center}
    Read off the counts and compute the overall error rate, the sensitivity, the specificity and
    the precision (positive predictive value) for the positive class \texttt{1}. Which single
    number would trace out the ROC curve if it were varied, and in which direction would
    sensitivity and specificity move?
  \item \pp{4} Consider $K$-nearest-neighbours. (i) Using the training set
    $(x,y)$: $(1,3)$, $(4,6)$, $(6,8)$, $(7,7)$, $(10,12)$, predict the response at $x_0=5$ with $k=3$
    (KNN regression, unweighted average). (ii) Explain how $k$ governs the bias--variance
    trade-off, which extreme ($k=1$ or $k$ large) is the most flexible, and why the
    \emph{training} error of $1$-NN is always $0$.
  \item \pp{4} Explain in one or two sentences what $k$-fold cross-validation does, then give
    one reason to prefer $k=5$ or $10$ over \emph{LOOCV} and one reason to prefer it over a
    \emph{single} validation split.
  \item \pp{2} You have $p=2{,}000$ predictors but expect only a handful to be truly related to
    the response. Ridge or lasso --- which penalty is more appropriate, and where would
    principal-components regression (PCR) sit relative to your choice?
  \item \pp{6} Match each scenario to exactly one method from
    \{\emph{linear regression}, \emph{lasso}, \emph{logistic regression},
    \emph{random forest}\} (each method used exactly once) and justify each choice in one
    sentence:
    \begin{enumerate}[label=(\roman*),nosep,topsep=1mm]
      \item Estimate how much one additional year of schooling adds to annual earnings, with a
        confidence interval, from $n=5{,}000$ survey responses.
      \item Predict a continuous blood-pressure outcome from $p=8{,}000$ genetic markers for
        $n=150$ patients, and report a short list of the relevant markers.
      \item Classify each incoming e-mail as spam or not, and explain to auditors how each
        feature changes the odds of an e-mail being spam.
      \item Predict machine failure from hundreds of mixed sensor variables with strong
        non-linear effects and interactions; only predictive accuracy matters.
    \end{enumerate}
\end{enumerate}

\ifdefined\withsolutions
\begin{solutionbox}
\textbf{(a) Reading the \texttt{confusion\_matrix} output.}
The array is ordered $[0,1]$ with rows $=$ true and columns $=$ predicted, so the diagonal holds
the correct predictions: from the table $TN=120$ and $FP=30$ (true \texttt{0}), $FN=10$ and
$TP=40$ (true \texttt{1}); total $200$, with $50$ diseased and $150$ healthy.
\[
  \text{error rate}=\frac{FP+FN}{200}=\frac{30+10}{200}=\mathbf{0.200},\qquad
  \text{sensitivity}=\frac{TP}{TP+FN}=\frac{40}{50}=\mathbf{0.800},
\]
\[
  \text{specificity}=\frac{TN}{TN+FP}=\frac{120}{150}=\mathbf{0.800},\qquad
  \text{precision}=\frac{TP}{TP+FP}=\frac{40}{70}=\mathbf{0.571}.
\]
Varying the \emph{classification threshold} on the predicted probability traces out the ROC
curve. Lowering the threshold labels more patients ``diseased'', which \emph{raises}
sensitivity (fewer missed cases) but \emph{lowers} specificity (more false alarms); raising it
does the reverse --- the ROC curve captures exactly this trade-off.

\medskip
\textbf{(b) $K$-nearest-neighbours.}
(i) Distances from $x_0=5$ are $|1-5|=4$, $|4-5|=1$, $|6-5|=1$, $|7-5|=2$, $|10-5|=5$; the three
nearest neighbours are $x=4,6,7$ with responses $6,8,7$, so
$\hat{y}(5)=(6+8+7)/3=21/3=\mathbf{7}$.
(ii) Small $k$ gives a very flexible fit: low bias but high variance, since the prediction hangs
on a few nearby points and a wiggly boundary. Large $k$ averages over many neighbours, giving a
smooth, stable fit: high bias but low variance. The most flexible extreme is $\mathbf{k=1}$
(flexibility $\propto 1/k$). The training error of $1$-NN is always $0$ because each training
point's single nearest neighbour is \emph{itself}, so it predicts its own observed response
exactly --- which is why training error is no guide to test performance.

\medskip
\textbf{(c) $k$-fold cross-validation.}
It splits the data into $k$ roughly equal folds; each fold in turn serves as the validation set
while the model is trained on the other $k-1$ folds, and the $k$ resulting errors are averaged.
\emph{vs.\ LOOCV:} the $n$ LOOCV fits use almost identical training sets, so their errors are
highly correlated and their average has \emph{higher variance} than $5$- or $10$-fold CV (and
LOOCV costs $n$ fits, absent shortcut formulas).
\emph{vs.\ a single validation split:} a single split gives a high-variance estimate that
depends heavily on the particular partition and is biased upward because only part of the data
trains the model; $k$-fold CV averages over $k$ folds and trains on $(k-1)/k$ of the data,
giving a more stable, less biased estimate.

\medskip
\textbf{(d) Lasso.}
The $\ell_1$ penalty drives many coefficients \emph{exactly to zero} and thus performs variable
selection --- ideal when the truth is sparse (a few real signals among $2{,}000$ predictors) and
an interpretable small model is wanted. Ridge only shrinks coefficients and keeps all $2{,}000$
predictors, so it tends to lose here. PCR is likewise \emph{not} sparse: it regresses on a few
principal components, so all original variables load on the retained components and no variable
is truly dropped --- like ridge it does not perform selection, so it does not deliver the short
interpretable list that lasso does.

\medskip
\textbf{(e) Matching.}
\begin{itemize}[nosep,leftmargin=5mm]
  \item (i) \textbf{Linear regression}: a quantitative response with $n\gg p$ and the goal of
    \emph{inference} --- an effect size with a confidence interval --- exactly what OLS with
    standard errors provides.
  \item (ii) \textbf{Lasso}: $p\gg n$ makes OLS infeasible; the $\ell_1$ penalty regularises the
    fit and returns a short list of relevant markers.
  \item (iii) \textbf{Logistic regression}: a binary response with interpretable coefficients
    ($e^{\hat\beta_j}$ are odds ratios) --- well suited to explaining odds to auditors.
  \item (iv) \textbf{Random forest}: automatically captures non-linearities and interactions,
    is accurate and robust with many mixed predictors, and interpretability is not required.
\end{itemize}
\end{solutionbox}
\fi

\vspace{6mm}
\begin{center}
\rule{0.35\textwidth}{0.4pt}\\[1mm]
\textit{End of the examination --- good luck!}
\end{center}

\end{document}
"
% ============================================================
\documentclass[11pt,a4paper]{article}
\usepackage[utf8]{inputenc}
\usepackage[T1]{fontenc}
\usepackage[english]{babel}
\usepackage[margin=2.2cm]{geometry}
\usepackage{amsmath,amssymb}
\usepackage{booktabs,array}
\usepackage{enumitem}
\usepackage{xcolor}
\usepackage{tcolorbox}
\tcbuselibrary{skins,breakable}

\setlength{\parindent}{0pt}

% Teal solution box (only shown when \withsolutions is defined)
\newtcolorbox{solutionbox}{enhanced,breakable,colback=teal!5,colframe=teal!55!black,
  boxrule=0pt,leftrule=2.5pt,arc=2pt,fonttitle=\bfseries,title=Solution,
  left=2mm,right=2mm,top=1mm,bottom=1.5mm}

\newcommand{\problem}[2]{\par\vspace{7mm}\noindent{\large\textbf{Problem #1}}%
  \hfill{\large\textbf{(#2 points)}}\par\nopagebreak\vspace{0.6mm}%
  \noindent\rule{\textwidth}{0.6pt}\par\nopagebreak\vspace{2.5mm}}
\newcommand{\pp}[1]{\textbf{[#1~P]}\enspace}

\begin{document}

% ------------------------------------------------------------
% Header block
% ------------------------------------------------------------
\begin{center}
  {\Large\textbf{HSBI --- Bielefeld University of Applied Sciences and Arts}}\\[1.2mm]
  {\large Quantitative Research Methods}\\[1mm]
  {\normalsize Summer Semester 2026 \quad$\cdot$\quad Examiner: Prof.\ Dr.\ Christoph Weisser}\\[4.5mm]
  {\LARGE\textbf{Final Mock Examination --- Set B (Lectures 1--12)}}\\[1.2mm]
  {\normalsize Comprehensive coverage of ISLP Chapters 2--8, 10 and 13,\\
   with emphasis on Chapters 7 (splines/GAMs), 8 (trees), 10 (deep learning) and 13 (multiple testing)}%
  \ifdefined\withsolutions\\[2.2mm]{\large\textcolor{teal!55!black}{\textbf{--- Version with detailed solutions ---}}}\fi
\end{center}

\vspace{2.5mm}
\begin{center}
\begin{tabular}{@{}ll@{}}
  \textbf{Duration:} & \textbf{120 minutes}\\
  \textbf{Total credit:} & \textbf{120 points}\\
  \textbf{Allowed aids:} & non-programmable calculator; one handwritten A4 sheet of notes (both sides)\\
\end{tabular}
\end{center}

\vspace{3mm}
Name: \rule[-1mm]{0.38\textwidth}{0.4pt}\hfill Matriculation no.: \rule[-1mm]{0.24\textwidth}{0.4pt}

\vspace{4.5mm}
\begin{center}
\renewcommand{\arraystretch}{1.35}
\begin{tabular}{|l|c|c|c|c|c|c|c|c|}
\hline
\textbf{Problem} & \textbf{P1} & \textbf{P2} & \textbf{P3} & \textbf{P4} & \textbf{P5} & \textbf{P6} & \textbf{P7} & \textbf{Total}\\\hline
Maximum points & 16 & 8 & 20 & 16 & 20 & 20 & 20 & \textbf{120}\\\hline
Achieved points & \hspace{8mm} & \hspace{8mm} & \hspace{8mm} & \hspace{8mm} & \hspace{8mm} & \hspace{8mm} & \hspace{8mm} & \hspace{10mm}\\\hline
\end{tabular}
\end{center}

\vspace{2.5mm}
\textbf{Instructions}
\begin{itemize}[nosep,leftmargin=5.5mm]
  \item Justify all answers. Numerical results without a derivation or a brief reasoning receive no credit.
  \item Round final numerical results to 3 decimal places unless stated otherwise. Small deviations caused by carrying rounded intermediate results are accepted.
  \item All numerical constants you need (powers of $e$, logarithms, $(0.95)^{10}$, $t$-quantiles) are provided in the problems.
  \item The number of points per problem roughly equals the recommended working time in minutes.
\end{itemize}
\vspace{1mm}\noindent\rule{\textwidth}{0.6pt}

% ============================================================
\problem{1: Moving beyond linearity --- splines}{16}
{\small\itshape Lecture slides: Chapter 7 --- Moving Beyond Linearity (Lecture 9).}\par\vspace{1mm}

Throughout this problem, consider regression splines built on a single predictor $X$ using
$K=5$ interior knots.

\begin{enumerate}[label=(\alph*),leftmargin=7mm,itemsep=2.5mm]
  \item \pp{4} A \emph{cubic spline} is a piecewise cubic polynomial that is continuous and has
    continuous first and second derivatives at every knot. How many degrees of freedom
    (free parameters, counting the intercept) does a cubic spline with $K$ interior knots use?
    Justify your answer by counting the parameters of the polynomial pieces and subtracting the
    constraints at the knots, and evaluate for $K=5$.
  \item \pp{4} How many degrees of freedom does a \emph{natural} cubic spline with the same
    $K=5$ knots use? State the additional constraint a natural spline imposes, in which region
    of the predictor space it acts, and why practitioners find this constraint useful.
  \item \pp{4} A \emph{smoothing spline} is the function $g$ that minimises
    \[
      \sum_{i=1}^{n}\bigl(y_i-g(x_i)\bigr)^2+\lambda\int g''(t)^2\,dt .
    \]
    Explain the meaning of the two terms, and describe the shape of the solution $\hat{g}$ in
    the two limiting cases $\lambda\to 0$ and $\lambda\to\infty$, giving the effective degrees
    of freedom that $\hat{g}$ approaches in each case.
  \item \pp{4} Write down the \emph{truncated-power basis} representation of a cubic spline with
    a single knot at $\xi$, define the truncated-power basis function $(x-\xi)^3_+$ explicitly,
    and explain in one sentence why adding this term keeps the fit continuous with continuous
    first and second derivatives at $\xi$.
\end{enumerate}

\ifdefined\withsolutions
\begin{solutionbox}
\textbf{(a) Degrees of freedom of a cubic spline.}
With $K$ interior knots the range of $X$ is divided into $K+1$ intervals, and on each interval
the fit is a cubic with $4$ coefficients, giving $4(K+1)=4K+4$ free parameters. At every knot
the requirement that $g$, $g'$ and $g''$ join continuously imposes $3$ constraints, i.e.\ $3K$
constraints in total. Hence
\[
  \mathrm{df}=4(K+1)-3K=K+4 .
\]
For $K=5$: $\mathrm{df}=5+4=\mathbf{9}$ (equivalently the basis intercept, $x$, $x^2$, $x^3$
plus one truncated-power term per knot, $4+K$).

\medskip
\textbf{(b) Natural cubic spline.}
A natural spline adds the requirement that the function be \emph{linear beyond the two boundary
knots}, i.e.\ on the outer regions $X<\xi_1$ and $X>\xi_K$. Forcing the quadratic and cubic
terms to vanish in each of those two regions removes $2$ parameters per side, $4$ in total:
\[
  \mathrm{df}=(K+4)-4=K,\qquad K=5:\ \mathrm{df}=\mathbf{5}.
\]
Usefulness: ordinary polynomials and unconstrained splines behave erratically near the
boundaries, where the data are sparse; forcing linearity there reduces the \emph{variance} of
predictions at and beyond the outer knots, at the cost of only a little extra bias.

\medskip
\textbf{(c) Smoothing spline.}
The first term (the RSS) rewards closeness to the observed data, while the second term
penalises \emph{roughness}: $g''$ measures the local curvature of $g$, so $\int g''(t)^2dt$ is
large for wiggly functions, and $\lambda\ge 0$ controls the trade-off.
\emph{$\lambda\to 0$:} the penalty vanishes, so $\hat{g}$ may be arbitrarily wiggly and
\emph{interpolates} every training point (zero training RSS, extreme overfitting); the
effective degrees of freedom approach $n$.
\emph{$\lambda\to\infty$:} any curvature becomes infinitely costly, forcing $g''\equiv 0$, so
$\hat{g}$ is a straight line, namely the ordinary \emph{least-squares regression line}
(effective df $\to 2$). For intermediate $\lambda$ (chosen e.g.\ by LOOCV) the effective df
decrease smoothly from $n$ down to $2$.

\medskip
\textbf{(d) Truncated-power basis, single knot $\xi$.}
\[
  f(x)=\beta_0+\beta_1x+\beta_2x^2+\beta_3x^3+\beta_4\,(x-\xi)^3_+,
  \qquad
  (x-\xi)^3_+=\begin{cases}(x-\xi)^3, & x>\xi\\[0.5mm] 0, & x\le\xi\end{cases}
\]
At $x=\xi$ the added function $(x-\xi)^3_+$ and its first and second derivatives are all $0$, so
including it alters only the \emph{third} derivative across the knot; the fit therefore stays
continuous with continuous first and second derivatives --- precisely the smoothness that
defines a cubic spline.
\end{solutionbox}
\fi

% ============================================================
\problem{2: Generalised additive models}{8}
{\small\itshape Lecture slides: Chapter 7 --- Moving Beyond Linearity (Lecture 9).}\par\vspace{1mm}

\begin{enumerate}[label=(\alph*),leftmargin=7mm,itemsep=2.5mm]
  \item \pp{3} Write down the general form of a \emph{generalised additive model} (GAM) for a
    quantitative response $Y$ and predictors $X_1,\dots,X_p$, and explain in one sentence how it
    generalises multiple linear regression.
  \item \pp{3} State one important \emph{interpretability advantage} of GAMs (concerning the
    fitted functions $f_j$) and their central structural \emph{limitation}, including how that
    limitation can be repaired.
  \item \pp{2} Describe in one or two sentences how the \emph{backfitting} algorithm fits a GAM.
\end{enumerate}

\ifdefined\withsolutions
\begin{solutionbox}
\textbf{(a) Model form.}
\[
  Y=\beta_0+f_1(X_1)+f_2(X_2)+\dots+f_p(X_p)+\varepsilon .
\]
Multiple linear regression is the special case in which every $f_j(X_j)=\beta_jX_j$; a GAM
replaces each straight-line term by a flexible smooth function $f_j$ (e.g.\ a natural or
smoothing spline, or a local regression) while retaining the \emph{additive} structure.

\medskip
\textbf{(b) Advantage and limitation.}
\emph{Advantage (interpretability):} because the effects add up, the contribution of each
predictor is carried by its own function $f_j$, which can be plotted and read off one panel at
a time with the other predictors held fixed --- non-linear effects stay as easy to inspect as
coefficients in a linear model.
\emph{Limitation:} pure additivity rules out \emph{interactions} --- the effect of $X_j$ is not
allowed to depend on the value of another $X_k$. Repair: introduce interaction terms by hand,
e.g.\ a product $X_j\times X_k$ or a low-dimensional smooth surface $f_{jk}(X_j,X_k)$ as an
extra additive component.

\medskip
\textbf{(c) Backfitting.}
Loop over the predictors and repeatedly re-estimate one $f_j$ at a time (with a smoother) using
the \emph{partial residuals} $r_i=y_i-\beta_0-\sum_{k\ne j}f_k(x_{ik})$, keeping the other
fitted functions fixed; repeat the cycles until the $f_j$ no longer change appreciably.
\end{solutionbox}
\fi

% ============================================================
\problem{3: Regression and classification trees by hand}{20}
{\small\itshape Lecture slides: Chapter 8 --- Tree-Based Methods (Lecture 10).}\par\vspace{1mm}

A logistics firm records, for $n=6$ regional depots, the number of delivery vans $x$ and the
number of parcels delivered per day $y$ (in thousands):

\begin{center}
\begin{tabular}{lcccccc}
\toprule
Depot $i$ & 1 & 2 & 3 & 4 & 5 & 6\\
\midrule
Vans $x_i$   & 10 & 12 & 14 & 16 & 18 & 20\\
Parcels $y_i$ & 6  & 4  & 5  & 11 & 13 & 12\\
\bottomrule
\end{tabular}
\end{center}

A regression tree with a \emph{single} split at cutpoint $s$ partitions the depots into
$R_1=\{i: x_i<s\}$ and $R_2=\{i: x_i\ge s\}$ and predicts the region mean $\hat{y}_{R_m}$ in
each region.

\begin{enumerate}[label=(\alph*),leftmargin=7mm,itemsep=2.5mm]
  \item \pp{8} For each of the two candidate cutpoints $s=15$ and $s=17$, determine the two
    regions, their means, and the total residual sum of squares
    $\operatorname{RSS}(s)=\sum_{m=1}^{2}\sum_{i\in R_m}(y_i-\hat{y}_{R_m})^2$.
    Which cutpoint does recursive binary splitting prefer, and why?
  \item \pp{2} Suppose recursive binary splitting selects the split at $s=15$, giving regions
    $R_1=\{x<15\}$ with mean $\hat{y}_{R_1}=5$ and $R_2=\{x\ge 15\}$ with mean
    $\hat{y}_{R_2}=12$. Predict the number of parcels for a new depot with $x=13$ vans.
  \item \pp{6} A \emph{classification} tree for a different task has a terminal node containing
    $12$ training observations from three classes: $6$ ``Premium'', $3$ ``Standard'' and $3$
    ``Basic''. Compute the class proportions, the Gini index
    $G=\sum_{k}\hat{p}_{k}(1-\hat{p}_{k})$ and the entropy $D=-\sum_{k}\hat{p}_{k}\ln\hat{p}_{k}$
    of this node \emph{(use $\ln 0.5=-0.6931$ and $\ln 0.25=-1.3863$)}. Which class does the node
    predict?
  \item \pp{4} In \emph{cost-complexity pruning} one minimises, over subtrees $T\subseteq T_0$
    of the large tree $T_0$,
    \[
      \sum_{m=1}^{|T|}\ \sum_{i:\,x_i\in R_m}\bigl(y_i-\hat{y}_{R_m}\bigr)^2+\alpha\,|T| .
    \]
    Explain the role of $\alpha$: what happens for $\alpha=0$, what happens as $\alpha$ grows,
    and how is $\alpha$ chosen in practice?
\end{enumerate}

\ifdefined\withsolutions
\begin{solutionbox}
\textbf{(a) Evaluating the two cutpoints.}
For reference the root mean is $\bar{y}=(6+4+5+11+13+12)/6=51/6=8.5$.

\emph{Cutpoint $s=15$:} $R_1=\{1,2,3\}$ with $y$-values $\{6,4,5\}$, mean $\hat{y}_{R_1}=15/3=5$;
$R_2=\{4,5,6\}$ with $y$-values $\{11,13,12\}$, mean $\hat{y}_{R_2}=36/3=12$.
\[
  \begin{aligned}
  \operatorname{RSS}_{R_1}&=(6-5)^2+(4-5)^2+(5-5)^2=1+1+0=2,\\
  \operatorname{RSS}_{R_2}&=(11-12)^2+(13-12)^2+(12-12)^2=1+1+0=2,
  \end{aligned}
\]
so $\operatorname{RSS}(15)=2+2=\mathbf{4}$.

\emph{Cutpoint $s=17$:} $R_1=\{1,2,3,4\}$ with $y$-values $\{6,4,5,11\}$, mean
$\hat{y}_{R_1}=26/4=6.5$; $R_2=\{5,6\}$ with $y$-values $\{13,12\}$, mean $\hat{y}_{R_2}=12.5$.
\[
  \operatorname{RSS}_{R_1}=(6-6.5)^2+(4-6.5)^2+(5-6.5)^2+(11-6.5)^2=0.25+6.25+2.25+20.25=29,
\]
\[
  \operatorname{RSS}_{R_2}=(13-12.5)^2+(12-12.5)^2=0.25+0.25=0.5,
\]
so $\operatorname{RSS}(17)=29+0.5=\mathbf{29.5}$.

Recursive binary splitting is \emph{greedy} and takes the cutpoint with the smaller RSS: since
$4<29.5$, it prefers $\mathbf{s=15}$, which cleanly separates the three low-throughput depots
(mean $5$) from the three high-throughput depots (mean $12$).

\medskip
\textbf{(b) Prediction for $x=13$.}
$x=13<15$, so the new depot falls into $R_1$ and the tree predicts
$\hat{y}=\hat{y}_{R_1}=\mathbf{5}$ (i.e.\ $5{,}000$ parcels per day).

\medskip
\textbf{(c) Node impurity (three classes).}
$\hat{p}_{\text{Prem}}=6/12=0.5$, $\hat{p}_{\text{Std}}=3/12=0.25$,
$\hat{p}_{\text{Basic}}=3/12=0.25$.
\[
  G=0.5(1-0.5)+0.25(1-0.25)+0.25(1-0.25)=0.25+0.1875+0.1875=\mathbf{0.625}.
\]
\[
  \begin{aligned}
  D&=-\bigl(0.5\ln 0.5+0.25\ln 0.25+0.25\ln 0.25\bigr)\\
   &=-\bigl(0.5(-0.6931)+0.5(-1.3863)\bigr)=0.3466+0.6931=\mathbf{1.040}.
  \end{aligned}
\]
The node predicts the majority class, \textbf{Premium} (estimated probability $0.5$). Both $G$
and $D$ would be $0$ for a pure node; here the node is quite impure, consistent with the split
carrying three classes.

\medskip
\textbf{(d) Role of $\alpha$.}
$\alpha\ge 0$ is the price charged per terminal node, trading the subtree's fit (first term)
against its size $|T|$.
For $\alpha=0$ the criterion reduces to the training RSS, so the full unpruned tree $T_0$ is
optimal. As $\alpha$ increases, each extra leaf must ``earn its keep''; a branch whose RSS
reduction is smaller than $\alpha$ is pruned away, so the optimal subtree shrinks and one
obtains a \emph{nested sequence} of subtrees as $\alpha$ climbs (weakest-link pruning). In
practice $\alpha$ is selected by \emph{cross-validation} --- pick the $\alpha$ giving the lowest
CV error, then refit that subtree on the full data --- which balances bias and variance rather
than merely minimising training error.
\end{solutionbox}
\fi

% ============================================================
\problem{4: Ensembles --- bagging, random forests, boosting}{16}
{\small\itshape Lecture slides: Chapter 8 --- Tree-Based Methods (Lecture 10).}\par\vspace{1mm}

\begin{enumerate}[label=(\alph*),leftmargin=7mm,itemsep=2.5mm]
  \item \pp{6} Let $\hat{f}_1(x),\dots,\hat{f}_B(x)$ be $B$ identically distributed (but not
    independent) tree predictions at a fixed point $x$, each with variance $\sigma^2$ and
    pairwise correlation $\rho$. The variance of the bagged average is
    \[
      \operatorname{Var}\Bigl(\tfrac{1}{B}\sum_{b=1}^{B}\hat{f}_b(x)\Bigr)
      =\rho\,\sigma^{2}+\frac{1-\rho}{B}\,\sigma^{2}.
    \]
    (i) Explain briefly why this formula shows that bagging reduces variance and which term
    limits the reduction. (ii) Evaluate the variance for $B=500$, $\rho=0.7$ and $\sigma^2=4$,
    and compare it with the variance of a single tree. (iii) What value does the variance
    approach as $B\to\infty$?
  \item \pp{4} How do \emph{random forests} modify bagging at each split, and what is the
    recommended subset size $m$ for \emph{classification} as a function of $p$ (evaluate for
    $p=36$)? Using the variance formula
    $\operatorname{Var}=\rho\sigma^2+\tfrac{1-\rho}{B}\sigma^2$ from part~(a), explain why this
    modification improves on bagging.
  \item \pp{4} A data scientist fits a boosting model for regression trees in
    \texttt{scikit-learn}:
    \begin{center}
    \texttt{GradientBoostingRegressor(n\_estimators=500, learning\_rate=0.01, max\_depth=1)}
    \end{center}
    These three arguments are exactly the tuning parameters of boosting. State in one sentence
    each what \texttt{n\_estimators}, \texttt{learning\_rate} and \texttt{max\_depth} control,
    and explain the ``slow learning'' idea behind the small \texttt{learning\_rate}.
  \item \pp{2} What does the \emph{out-of-bag} (OOB) error of a bagged model estimate, and why
    is it essentially free to obtain?
\end{enumerate}

\ifdefined\withsolutions
\begin{solutionbox}
\textbf{(a) Variance of the bagged average.}
(i) The second term $\frac{1-\rho}{B}\sigma^2$ shrinks to $0$ as $B$ grows, so averaging over
many trees removes the ``independent part'' of their variability --- this is the variance
reduction of bagging. The first term $\rho\sigma^2$ does \emph{not} depend on $B$: the
bootstrap trees are grown on overlapping resamples of the same data (and share the same strong
predictors), so they are positively correlated, and this common part of the variance cannot be
averaged away. It is the floor that limits the reduction.
(ii) For $B=500$, $\rho=0.7$, $\sigma^2=4$:
\[
  \operatorname{Var}=0.7\cdot 4+\frac{1-0.7}{500}\cdot 4
  =2.8+\frac{0.3\cdot 4}{500}=2.8+0.0024=\mathbf{2.802},
\]
versus $\sigma^2=4$ for a single tree --- a reduction of about $30\%$.
(iii) As $B\to\infty$ the variance approaches the floor $\rho\sigma^2=0.7\cdot 4=\mathbf{2.8}$.
Enlarging $B$ further does not overfit; it only stabilises the average.

\medskip
\textbf{(b) Random forests.}
At each split the tree may consider only a fresh \emph{random subset} of $m$ out of the $p$
predictors as candidates. For classification the usual choice is $m\approx\sqrt{p}$, so for
$p=36$: $m=\sqrt{36}=\mathbf{6}$.
Why it helps: in plain bagging a dominant predictor is chosen at the top split of nearly every
tree, making the trees look alike and pushing $\rho$ up. Withholding that predictor from about
$(p-m)/p$ of the split searches \emph{decorrelates} the trees, i.e.\ lowers $\rho$; by (a) this
lowers the irreducible floor $\rho\sigma^2$ of the ensemble variance, which plain averaging can
never reach.

\medskip
\textbf{(c) Boosting tuning parameters (read from the call).}
\begin{itemize}[nosep,leftmargin=5mm]
  \item \texttt{n\_estimators}$=B=500$: the number of small trees added in sequence; unlike
    bagging, boosting \emph{can} overfit if $B$ is too large, so $B$ is tuned by cross-validation.
  \item \texttt{learning\_rate}$=\lambda=0.01$ (shrinkage): scales the contribution of each new
    tree and hence the speed of learning; a smaller $\lambda$ usually requires a larger $B$.
  \item \texttt{max\_depth}$=d=1$ (interaction depth): the number of splits per tree and the
    interaction order captured; $d=1$ means \emph{stumps} (an additive model).
\end{itemize}
\emph{Slow learning:} each new tree is fit to the current \emph{residuals}, and only the small
fraction \texttt{learning\_rate}$=\lambda$ of it is added to the model, so the fit improves
gradually and precisely where it is still weak, rather than fitting the data hard in a single
step --- which tends to generalise better.

\medskip
\textbf{(d) OOB error.}
Each bootstrap sample omits, on average, about one third of the observations
(``out-of-bag''). Predicting each observation only from the trees that did \emph{not} train on
it, and aggregating, gives the OOB error --- a valid estimate of the bagged model's
\emph{test} error. It is essentially free because it reuses a by-product of the bootstrap; no
separate validation set or cross-validation loop is needed.
\end{solutionbox}
\fi

% ============================================================
\problem{5: Neural networks by hand}{20}
{\small\itshape Lecture slides: Chapter 10 --- Deep Learning (Lecture 11).}\par\vspace{1mm}

Consider a small feed-forward network for a quantitative response with input
$x=(x_1,x_2)=(2,1)$, one hidden layer of two ReLU units ($g(z)=\max\{0,z\}$) and a linear
output unit:
\[
  A_1=g\bigl(-2+3\,x_1-1\,x_2\bigr),\qquad
  A_2=g\bigl(1-2\,x_1+1\,x_2\bigr),\qquad
  f(x)=-1+2\,A_1+3\,A_2 .
\]

\begin{enumerate}[label=(\alph*),leftmargin=7mm,itemsep=2.5mm]
  \item \pp{8} Carry out the forward pass step by step: compute both pre-activations, both
    activations $A_1,A_2$ and the output $f(x)$. Comment briefly on what the ReLU did to the
    second unit and why non-linear activations are essential.
  \item \pp{6} A classification network with $K=3$ classes returns the logits (output-layer
    pre-activations) $z=(z_1,z_2,z_3)=(2.0,\ 1.0,\ -1.0)$ for one observation whose true class
    is class~1. Compute all three softmax probabilities and the cross-entropy loss
    $-\ln\hat{p}_{\text{true}}$.
    \emph{(Use $e^{2}=7.389$, $e^{1}=2.718$, $e^{-1}=0.368$ and $\ln 10.475=2.349$.)}
  \item \pp{4} Count the parameters of a fully connected network with $p=6$ inputs, one hidden
    layer of $k=4$ units and an output layer of $K=3$ units, where every hidden and output unit
    carries a bias. Show the count separately for the two weight layers.
  \item \pp{2} A convolutional layer receives a $32\times 32$ image and applies filters of size
    $F=5$ with padding $P=0$ and stride $S=1$. Compute the spatial size of the output feature
    map from $\bigl(W-F+2P\bigr)/S+1$.
\end{enumerate}

\ifdefined\withsolutions
\begin{solutionbox}
\textbf{(a) Forward pass.}
Pre-activations:
\[
  z_1=-2+3\cdot 2-1\cdot 1=-2+6-1=3,\qquad
  z_2=1-2\cdot 2+1\cdot 1=1-4+1=-2 .
\]
ReLU activations:
\[
  A_1=g(3)=\max\{0,3\}=3,\qquad A_2=g(-2)=\max\{0,-2\}=0 .
\]
Output:
\[
  f(x)=-1+2\cdot 3+3\cdot 0=-1+6+0=\mathbf{5}.
\]
The ReLU \emph{switched the second unit off}: its negative pre-activation is clipped to $0$, so
unit~2 contributes nothing for this particular input (though it could for others). This
thresholding is exactly the source of non-linearity: with a purely linear $g$, a network of any
depth would collapse to a single linear function of $x$.

\medskip
\textbf{(b) Softmax and cross-entropy.}
Exponentials $e^{2.0}=7.389$, $e^{1.0}=2.718$, $e^{-1.0}=0.368$ sum to
$7.389+2.718+0.368=10.475$. Softmax probabilities:
\[
  \hat{p}_1=\frac{7.389}{10.475}=0.705,\qquad
  \hat{p}_2=\frac{2.718}{10.475}=0.259,\qquad
  \hat{p}_3=\frac{0.368}{10.475}=0.035,
\]
which sum to $0.999\approx 1$. Cross-entropy loss for the true class~1:
\[
  -\ln\hat{p}_1=-\ln\frac{e^{2}}{10.475}=-\bigl(2-\ln 10.475\bigr)=2.349-2=\mathbf{0.349}.
\]
(Equivalently $-\ln 0.705\approx 0.349$; a perfect prediction $\hat{p}_1=1$ would give loss 0.)

\medskip
\textbf{(c) Parameter count for $6\to 4\to 3$.}
Input $\to$ hidden: $6\cdot 4=24$ weights $+\,4$ biases $=28$.
Hidden $\to$ output: $4\cdot 3=12$ weights $+\,3$ biases $=15$.
Total: $28+15=\mathbf{43}$ parameters.

\medskip
\textbf{(d) Convolution output size.}
\[
  \frac{W-F+2P}{S}+1=\frac{32-5+2\cdot 0}{1}+1=\frac{27}{1}+1=28,
\]
so the output feature map is $\mathbf{28\times 28}$. With no padding a $5\times 5$ filter at
stride $1$ shrinks each spatial dimension by $F-1=4$.
\end{solutionbox}
\fi

% ============================================================
\problem{6: Multiple testing}{20}
{\small\itshape Lecture slides: Chapter 13 --- Multiple Testing (Lecture 12).}\par\vspace{1mm}

\begin{enumerate}[label=(\alph*),leftmargin=7mm,itemsep=2.5mm]
  \item \pp{4} A lab runs $m=10$ \emph{independent} hypothesis tests, each at level
    $\alpha=0.05$, and suppose all $10$ null hypotheses are true. Define the family-wise error
    rate (FWER) and compute it here \emph{(use $(0.95)^{10}=0.5987$)}. Interpret the result in
    one sentence.
  \item \pp{7} A team screens $m=6$ candidate biomarkers; the ordered $p$-values are
    \begin{center}
    \begin{tabular}{lcccccc}
    \toprule
    $j$ & 1 & 2 & 3 & 4 & 5 & 6\\
    \midrule
    $p_{(j)}$ & 0.002 & 0.009 & 0.012 & 0.015 & 0.040 & 0.300\\
    \bottomrule
    \end{tabular}
    \end{center}
    Control the FWER at $\alpha=0.05$ (i) with \emph{Bonferroni} and (ii) with the
    \emph{Holm step-down} procedure. Give the relevant thresholds, state exactly which
    hypotheses each method rejects, and explain why Holm can never reject fewer hypotheses than
    Bonferroni.
  \item \pp{6} Apply the \emph{Benjamini--Hochberg} (BH) procedure at FDR level $q=0.05$ to the
    same six ordered $p$-values ($m=6$; $p_{(1)},\dots,p_{(6)}=0.002,\ 0.009,\ 0.012,\ 0.015,\
    0.040,\ 0.300$): give the thresholds $jq/m$, determine which hypotheses are rejected, and
    state precisely what ``FDR control at $q=0.05$'' guarantees for the resulting set of
    rejections.
  \item \pp{3} A pharmaceutical company runs a \emph{confirmatory} Phase-III trial with $m=4$
    pre-specified primary endpoints; a single significant endpoint could support regulatory
    approval of the drug. Should the company control the FWER or the FDR? Justify your
    recommendation.
\end{enumerate}

\ifdefined\withsolutions
\begin{solutionbox}
\textbf{(a) FWER for $m=10$ independent tests.}
$\text{FWER}=\Pr(\text{at least one Type I error})=\Pr(V\ge 1)$, where $V$ counts the true
nulls that are falsely rejected. With independence and all nulls true,
\[
  \text{FWER}=1-\Pr(\text{no false rejection})=1-(1-0.05)^{10}=1-(0.95)^{10}
  =1-0.5987=\mathbf{0.4013}.
\]
Interpretation: although each individual test uses the $5\%$ level, there is about a $40\%$
chance of at least one false discovery across the ten tests --- testing many hypotheses without
adjustment makes spurious findings very likely.

\medskip
\textbf{(b) Bonferroni and Holm at $\alpha=0.05$.}
\emph{Bonferroni:} reject $H_{(j)}$ iff $p_{(j)}\le\alpha/m=0.05/6=0.00833$.
Only $p_{(1)}=0.002\le 0.00833$ \checkmark; already $p_{(2)}=0.009>0.00833$ $\times$.
$\Rightarrow$ Bonferroni rejects $H_{(1)}$ alone: \textbf{1 rejection}. (Note $p_{(2)}=0.009$ is
below the naive $0.05$ cutoff, yet Bonferroni does not reject it.)

\emph{Holm:} compare $p_{(j)}$ with $\alpha/(m-j+1)$ in order and stop at the first failure:
\begin{center}
\begin{tabular}{lcccccc}
\toprule
$j$ & 1 & 2 & 3 & 4 & 5 & 6\\
\midrule
$p_{(j)}$ & 0.002 & 0.009 & 0.012 & 0.015 & 0.040 & 0.300\\
$\alpha/(m-j+1)$ & 0.00833 & 0.0100 & 0.0125 & 0.0167 & 0.0250 & 0.0500\\
reject? & \checkmark & \checkmark & \checkmark & \checkmark & stop & ---\\
\bottomrule
\end{tabular}
\end{center}
$0.002\le 0.00833$, $0.009\le 0.0100$, $0.012\le 0.0125$, $0.015\le 0.0167$, but
$0.040>0.0250$: stop. $\Rightarrow$ Holm rejects $H_{(1)},H_{(2)},H_{(3)},H_{(4)}$:
\textbf{4 rejections}.

\emph{Why Holm $\supseteq$ Bonferroni:} Holm's thresholds $\alpha/(m-j+1)$ are all $\ge\alpha/m$
(with equality only at $j=1$), so any $p$-value that clears the strict Bonferroni bar also
clears Holm's more lenient later hurdles --- Holm is uniformly at least as powerful, while still
controlling the FWER at $\alpha$ with no independence assumption.

\medskip
\textbf{(c) Benjamini--Hochberg at $q=0.05$.}
Thresholds $jq/m=j\cdot 0.05/6$:
\begin{center}
\begin{tabular}{lcccccc}
\toprule
$j$ & 1 & 2 & 3 & 4 & 5 & 6\\
\midrule
$p_{(j)}$ & 0.002 & 0.009 & 0.012 & 0.015 & 0.040 & 0.300\\
$jq/m$ & 0.00833 & 0.0167 & 0.0250 & 0.0333 & 0.0417 & 0.0500\\
$p_{(j)}\le jq/m$? & yes & yes & yes & yes & \textbf{yes} & no\\
\bottomrule
\end{tabular}
\end{center}
$L=$ largest $j$ with $p_{(j)}\le jq/m$: here $j=5$ (since $0.040\le 0.0417$, while
$0.300>0.0500$). BH rejects \emph{all} $H_{(j)}$ with $j\le L$, i.e.\
$H_{(1)},\dots,H_{(5)}$: \textbf{5 rejections}.
Guarantee: $\text{FDR}=E\bigl[V/\max(R,1)\bigr]\le q=0.05$ --- \emph{on average} at most $5\%$
of the rejected hypotheses are false discoveries. It does \emph{not} promise that at most $5\%$
of these particular $5$ rejections are false, nor does it bound the probability of making any
false rejection (that would be FWER control).

\medskip
\textbf{(d) Confirmatory trial recommendation.}
Control the \textbf{FWER}. Here even a \emph{single} false positive is very costly --- it could
lead to approving an ineffective drug --- so the study must guard against any Type I error, not
merely keep their average proportion small. With only $m=4$ pre-specified endpoints, an FWER
procedure (e.g.\ Holm or Bonferroni at $0.05/4=0.0125$) sacrifices little power. FDR control,
appropriate for large exploratory screens where a few false leads are tolerable, is too
permissive for a confirmatory, decision-driving trial.
\end{solutionbox}
\fi

% ============================================================
\problem{7: Cumulative essentials (Chapters 2--6)}{20}
{\small\itshape Lecture slides: Chapters 2--6 (Lectures 1--8).}\par\vspace{1mm}

\begin{enumerate}[label=(\alph*),leftmargin=7mm,itemsep=2.5mm]
  \item \pp{4} A medical screening classifier is evaluated on $200$ patients (label
    \texttt{1}$=$diseased, \texttt{0}$=$healthy). In \texttt{scikit-learn},
    \texttt{confusion\_matrix(y\_true, y\_pred)} returns the array below (rows $=$ true class,
    columns $=$ predicted class; label order $[0,1]$):
    \begin{center}
    \begin{tabular}{lcc}
    \toprule
     & predicted \texttt{0} & predicted \texttt{1}\\
    \midrule
    true \texttt{0} (healthy)  & $120$ & $30$\\
    true \texttt{1} (diseased) & $10$  & $40$\\
    \bottomrule
    \end{tabular}
    \end{center}
    Read off the counts and compute the overall error rate, the sensitivity, the specificity and
    the precision (positive predictive value) for the positive class \texttt{1}. Which single
    number would trace out the ROC curve if it were varied, and in which direction would
    sensitivity and specificity move?
  \item \pp{4} Consider $K$-nearest-neighbours. (i) Using the training set
    $(x,y)$: $(1,3)$, $(4,6)$, $(6,8)$, $(7,7)$, $(10,12)$, predict the response at $x_0=5$ with $k=3$
    (KNN regression, unweighted average). (ii) Explain how $k$ governs the bias--variance
    trade-off, which extreme ($k=1$ or $k$ large) is the most flexible, and why the
    \emph{training} error of $1$-NN is always $0$.
  \item \pp{4} Explain in one or two sentences what $k$-fold cross-validation does, then give
    one reason to prefer $k=5$ or $10$ over \emph{LOOCV} and one reason to prefer it over a
    \emph{single} validation split.
  \item \pp{2} You have $p=2{,}000$ predictors but expect only a handful to be truly related to
    the response. Ridge or lasso --- which penalty is more appropriate, and where would
    principal-components regression (PCR) sit relative to your choice?
  \item \pp{6} Match each scenario to exactly one method from
    \{\emph{linear regression}, \emph{lasso}, \emph{logistic regression},
    \emph{random forest}\} (each method used exactly once) and justify each choice in one
    sentence:
    \begin{enumerate}[label=(\roman*),nosep,topsep=1mm]
      \item Estimate how much one additional year of schooling adds to annual earnings, with a
        confidence interval, from $n=5{,}000$ survey responses.
      \item Predict a continuous blood-pressure outcome from $p=8{,}000$ genetic markers for
        $n=150$ patients, and report a short list of the relevant markers.
      \item Classify each incoming e-mail as spam or not, and explain to auditors how each
        feature changes the odds of an e-mail being spam.
      \item Predict machine failure from hundreds of mixed sensor variables with strong
        non-linear effects and interactions; only predictive accuracy matters.
    \end{enumerate}
\end{enumerate}

\ifdefined\withsolutions
\begin{solutionbox}
\textbf{(a) Reading the \texttt{confusion\_matrix} output.}
The array is ordered $[0,1]$ with rows $=$ true and columns $=$ predicted, so the diagonal holds
the correct predictions: from the table $TN=120$ and $FP=30$ (true \texttt{0}), $FN=10$ and
$TP=40$ (true \texttt{1}); total $200$, with $50$ diseased and $150$ healthy.
\[
  \text{error rate}=\frac{FP+FN}{200}=\frac{30+10}{200}=\mathbf{0.200},\qquad
  \text{sensitivity}=\frac{TP}{TP+FN}=\frac{40}{50}=\mathbf{0.800},
\]
\[
  \text{specificity}=\frac{TN}{TN+FP}=\frac{120}{150}=\mathbf{0.800},\qquad
  \text{precision}=\frac{TP}{TP+FP}=\frac{40}{70}=\mathbf{0.571}.
\]
Varying the \emph{classification threshold} on the predicted probability traces out the ROC
curve. Lowering the threshold labels more patients ``diseased'', which \emph{raises}
sensitivity (fewer missed cases) but \emph{lowers} specificity (more false alarms); raising it
does the reverse --- the ROC curve captures exactly this trade-off.

\medskip
\textbf{(b) $K$-nearest-neighbours.}
(i) Distances from $x_0=5$ are $|1-5|=4$, $|4-5|=1$, $|6-5|=1$, $|7-5|=2$, $|10-5|=5$; the three
nearest neighbours are $x=4,6,7$ with responses $6,8,7$, so
$\hat{y}(5)=(6+8+7)/3=21/3=\mathbf{7}$.
(ii) Small $k$ gives a very flexible fit: low bias but high variance, since the prediction hangs
on a few nearby points and a wiggly boundary. Large $k$ averages over many neighbours, giving a
smooth, stable fit: high bias but low variance. The most flexible extreme is $\mathbf{k=1}$
(flexibility $\propto 1/k$). The training error of $1$-NN is always $0$ because each training
point's single nearest neighbour is \emph{itself}, so it predicts its own observed response
exactly --- which is why training error is no guide to test performance.

\medskip
\textbf{(c) $k$-fold cross-validation.}
It splits the data into $k$ roughly equal folds; each fold in turn serves as the validation set
while the model is trained on the other $k-1$ folds, and the $k$ resulting errors are averaged.
\emph{vs.\ LOOCV:} the $n$ LOOCV fits use almost identical training sets, so their errors are
highly correlated and their average has \emph{higher variance} than $5$- or $10$-fold CV (and
LOOCV costs $n$ fits, absent shortcut formulas).
\emph{vs.\ a single validation split:} a single split gives a high-variance estimate that
depends heavily on the particular partition and is biased upward because only part of the data
trains the model; $k$-fold CV averages over $k$ folds and trains on $(k-1)/k$ of the data,
giving a more stable, less biased estimate.

\medskip
\textbf{(d) Lasso.}
The $\ell_1$ penalty drives many coefficients \emph{exactly to zero} and thus performs variable
selection --- ideal when the truth is sparse (a few real signals among $2{,}000$ predictors) and
an interpretable small model is wanted. Ridge only shrinks coefficients and keeps all $2{,}000$
predictors, so it tends to lose here. PCR is likewise \emph{not} sparse: it regresses on a few
principal components, so all original variables load on the retained components and no variable
is truly dropped --- like ridge it does not perform selection, so it does not deliver the short
interpretable list that lasso does.

\medskip
\textbf{(e) Matching.}
\begin{itemize}[nosep,leftmargin=5mm]
  \item (i) \textbf{Linear regression}: a quantitative response with $n\gg p$ and the goal of
    \emph{inference} --- an effect size with a confidence interval --- exactly what OLS with
    standard errors provides.
  \item (ii) \textbf{Lasso}: $p\gg n$ makes OLS infeasible; the $\ell_1$ penalty regularises the
    fit and returns a short list of relevant markers.
  \item (iii) \textbf{Logistic regression}: a binary response with interpretable coefficients
    ($e^{\hat\beta_j}$ are odds ratios) --- well suited to explaining odds to auditors.
  \item (iv) \textbf{Random forest}: automatically captures non-linearities and interactions,
    is accurate and robust with many mixed predictors, and interpretability is not required.
\end{itemize}
\end{solutionbox}
\fi

\vspace{6mm}
\begin{center}
\rule{0.35\textwidth}{0.4pt}\\[1mm]
\textit{End of the examination --- good luck!}
\end{center}

\end{document}
"
% ============================================================
\documentclass[11pt,a4paper]{article}
\usepackage[utf8]{inputenc}
\usepackage[T1]{fontenc}
\usepackage[english]{babel}
\usepackage[margin=2.2cm]{geometry}
\usepackage{amsmath,amssymb}
\usepackage{booktabs,array}
\usepackage{enumitem}
\usepackage{xcolor}
\usepackage{tcolorbox}
\tcbuselibrary{skins,breakable}

\setlength{\parindent}{0pt}

% Teal solution box (only shown when \withsolutions is defined)
\newtcolorbox{solutionbox}{enhanced,breakable,colback=teal!5,colframe=teal!55!black,
  boxrule=0pt,leftrule=2.5pt,arc=2pt,fonttitle=\bfseries,title=Solution,
  left=2mm,right=2mm,top=1mm,bottom=1.5mm}

\newcommand{\problem}[2]{\par\vspace{7mm}\noindent{\large\textbf{Problem #1}}%
  \hfill{\large\textbf{(#2 points)}}\par\nopagebreak\vspace{0.6mm}%
  \noindent\rule{\textwidth}{0.6pt}\par\nopagebreak\vspace{2.5mm}}
\newcommand{\pp}[1]{\textbf{[#1~P]}\enspace}

\begin{document}

% ------------------------------------------------------------
% Header block
% ------------------------------------------------------------
\begin{center}
  {\Large\textbf{HSBI --- Bielefeld University of Applied Sciences and Arts}}\\[1.2mm]
  {\large Quantitative Research Methods}\\[1mm]
  {\normalsize Summer Semester 2026 \quad$\cdot$\quad Examiner: Prof.\ Dr.\ Christoph Weisser}\\[4.5mm]
  {\LARGE\textbf{Final Mock Examination --- Set B (Lectures 1--12)}}\\[1.2mm]
  {\normalsize Comprehensive coverage of ISLP Chapters 2--8, 10 and 13,\\
   with emphasis on Chapters 7 (splines/GAMs), 8 (trees), 10 (deep learning) and 13 (multiple testing)}%
  \ifdefined\withsolutions\\[2.2mm]{\large\textcolor{teal!55!black}{\textbf{--- Version with detailed solutions ---}}}\fi
\end{center}

\vspace{2.5mm}
\begin{center}
\begin{tabular}{@{}ll@{}}
  \textbf{Duration:} & \textbf{120 minutes}\\
  \textbf{Total credit:} & \textbf{120 points}\\
  \textbf{Allowed aids:} & non-programmable calculator; one handwritten A4 sheet of notes (both sides)\\
\end{tabular}
\end{center}

\vspace{3mm}
Name: \rule[-1mm]{0.38\textwidth}{0.4pt}\hfill Matriculation no.: \rule[-1mm]{0.24\textwidth}{0.4pt}

\vspace{4.5mm}
\begin{center}
\renewcommand{\arraystretch}{1.35}
\begin{tabular}{|l|c|c|c|c|c|c|c|c|}
\hline
\textbf{Problem} & \textbf{P1} & \textbf{P2} & \textbf{P3} & \textbf{P4} & \textbf{P5} & \textbf{P6} & \textbf{P7} & \textbf{Total}\\\hline
Maximum points & 16 & 8 & 20 & 16 & 20 & 20 & 20 & \textbf{120}\\\hline
Achieved points & \hspace{8mm} & \hspace{8mm} & \hspace{8mm} & \hspace{8mm} & \hspace{8mm} & \hspace{8mm} & \hspace{8mm} & \hspace{10mm}\\\hline
\end{tabular}
\end{center}

\vspace{2.5mm}
\textbf{Instructions}
\begin{itemize}[nosep,leftmargin=5.5mm]
  \item Justify all answers. Numerical results without a derivation or a brief reasoning receive no credit.
  \item Round final numerical results to 3 decimal places unless stated otherwise. Small deviations caused by carrying rounded intermediate results are accepted.
  \item All numerical constants you need (powers of $e$, logarithms, $(0.95)^{10}$, $t$-quantiles) are provided in the problems.
  \item The number of points per problem roughly equals the recommended working time in minutes.
\end{itemize}
\vspace{1mm}\noindent\rule{\textwidth}{0.6pt}

% ============================================================
\problem{1: Moving beyond linearity --- splines}{16}
{\small\itshape Lecture slides: Chapter 7 --- Moving Beyond Linearity (Lecture 9).}\par\vspace{1mm}

Throughout this problem, consider regression splines built on a single predictor $X$ using
$K=5$ interior knots.

\begin{enumerate}[label=(\alph*),leftmargin=7mm,itemsep=2.5mm]
  \item \pp{4} A \emph{cubic spline} is a piecewise cubic polynomial that is continuous and has
    continuous first and second derivatives at every knot. How many degrees of freedom
    (free parameters, counting the intercept) does a cubic spline with $K$ interior knots use?
    Justify your answer by counting the parameters of the polynomial pieces and subtracting the
    constraints at the knots, and evaluate for $K=5$.
  \item \pp{4} How many degrees of freedom does a \emph{natural} cubic spline with the same
    $K=5$ knots use? State the additional constraint a natural spline imposes, in which region
    of the predictor space it acts, and why practitioners find this constraint useful.
  \item \pp{4} A \emph{smoothing spline} is the function $g$ that minimises
    \[
      \sum_{i=1}^{n}\bigl(y_i-g(x_i)\bigr)^2+\lambda\int g''(t)^2\,dt .
    \]
    Explain the meaning of the two terms, and describe the shape of the solution $\hat{g}$ in
    the two limiting cases $\lambda\to 0$ and $\lambda\to\infty$, giving the effective degrees
    of freedom that $\hat{g}$ approaches in each case.
  \item \pp{4} Write down the \emph{truncated-power basis} representation of a cubic spline with
    a single knot at $\xi$, define the truncated-power basis function $(x-\xi)^3_+$ explicitly,
    and explain in one sentence why adding this term keeps the fit continuous with continuous
    first and second derivatives at $\xi$.
\end{enumerate}

\ifdefined\withsolutions
\begin{solutionbox}
\textbf{(a) Degrees of freedom of a cubic spline.}
With $K$ interior knots the range of $X$ is divided into $K+1$ intervals, and on each interval
the fit is a cubic with $4$ coefficients, giving $4(K+1)=4K+4$ free parameters. At every knot
the requirement that $g$, $g'$ and $g''$ join continuously imposes $3$ constraints, i.e.\ $3K$
constraints in total. Hence
\[
  \mathrm{df}=4(K+1)-3K=K+4 .
\]
For $K=5$: $\mathrm{df}=5+4=\mathbf{9}$ (equivalently the basis intercept, $x$, $x^2$, $x^3$
plus one truncated-power term per knot, $4+K$).

\medskip
\textbf{(b) Natural cubic spline.}
A natural spline adds the requirement that the function be \emph{linear beyond the two boundary
knots}, i.e.\ on the outer regions $X<\xi_1$ and $X>\xi_K$. Forcing the quadratic and cubic
terms to vanish in each of those two regions removes $2$ parameters per side, $4$ in total:
\[
  \mathrm{df}=(K+4)-4=K,\qquad K=5:\ \mathrm{df}=\mathbf{5}.
\]
Usefulness: ordinary polynomials and unconstrained splines behave erratically near the
boundaries, where the data are sparse; forcing linearity there reduces the \emph{variance} of
predictions at and beyond the outer knots, at the cost of only a little extra bias.

\medskip
\textbf{(c) Smoothing spline.}
The first term (the RSS) rewards closeness to the observed data, while the second term
penalises \emph{roughness}: $g''$ measures the local curvature of $g$, so $\int g''(t)^2dt$ is
large for wiggly functions, and $\lambda\ge 0$ controls the trade-off.
\emph{$\lambda\to 0$:} the penalty vanishes, so $\hat{g}$ may be arbitrarily wiggly and
\emph{interpolates} every training point (zero training RSS, extreme overfitting); the
effective degrees of freedom approach $n$.
\emph{$\lambda\to\infty$:} any curvature becomes infinitely costly, forcing $g''\equiv 0$, so
$\hat{g}$ is a straight line, namely the ordinary \emph{least-squares regression line}
(effective df $\to 2$). For intermediate $\lambda$ (chosen e.g.\ by LOOCV) the effective df
decrease smoothly from $n$ down to $2$.

\medskip
\textbf{(d) Truncated-power basis, single knot $\xi$.}
\[
  f(x)=\beta_0+\beta_1x+\beta_2x^2+\beta_3x^3+\beta_4\,(x-\xi)^3_+,
  \qquad
  (x-\xi)^3_+=\begin{cases}(x-\xi)^3, & x>\xi\\[0.5mm] 0, & x\le\xi\end{cases}
\]
At $x=\xi$ the added function $(x-\xi)^3_+$ and its first and second derivatives are all $0$, so
including it alters only the \emph{third} derivative across the knot; the fit therefore stays
continuous with continuous first and second derivatives --- precisely the smoothness that
defines a cubic spline.
\end{solutionbox}
\fi

% ============================================================
\problem{2: Generalised additive models}{8}
{\small\itshape Lecture slides: Chapter 7 --- Moving Beyond Linearity (Lecture 9).}\par\vspace{1mm}

\begin{enumerate}[label=(\alph*),leftmargin=7mm,itemsep=2.5mm]
  \item \pp{3} Write down the general form of a \emph{generalised additive model} (GAM) for a
    quantitative response $Y$ and predictors $X_1,\dots,X_p$, and explain in one sentence how it
    generalises multiple linear regression.
  \item \pp{3} State one important \emph{interpretability advantage} of GAMs (concerning the
    fitted functions $f_j$) and their central structural \emph{limitation}, including how that
    limitation can be repaired.
  \item \pp{2} Describe in one or two sentences how the \emph{backfitting} algorithm fits a GAM.
\end{enumerate}

\ifdefined\withsolutions
\begin{solutionbox}
\textbf{(a) Model form.}
\[
  Y=\beta_0+f_1(X_1)+f_2(X_2)+\dots+f_p(X_p)+\varepsilon .
\]
Multiple linear regression is the special case in which every $f_j(X_j)=\beta_jX_j$; a GAM
replaces each straight-line term by a flexible smooth function $f_j$ (e.g.\ a natural or
smoothing spline, or a local regression) while retaining the \emph{additive} structure.

\medskip
\textbf{(b) Advantage and limitation.}
\emph{Advantage (interpretability):} because the effects add up, the contribution of each
predictor is carried by its own function $f_j$, which can be plotted and read off one panel at
a time with the other predictors held fixed --- non-linear effects stay as easy to inspect as
coefficients in a linear model.
\emph{Limitation:} pure additivity rules out \emph{interactions} --- the effect of $X_j$ is not
allowed to depend on the value of another $X_k$. Repair: introduce interaction terms by hand,
e.g.\ a product $X_j\times X_k$ or a low-dimensional smooth surface $f_{jk}(X_j,X_k)$ as an
extra additive component.

\medskip
\textbf{(c) Backfitting.}
Loop over the predictors and repeatedly re-estimate one $f_j$ at a time (with a smoother) using
the \emph{partial residuals} $r_i=y_i-\beta_0-\sum_{k\ne j}f_k(x_{ik})$, keeping the other
fitted functions fixed; repeat the cycles until the $f_j$ no longer change appreciably.
\end{solutionbox}
\fi

% ============================================================
\problem{3: Regression and classification trees by hand}{20}
{\small\itshape Lecture slides: Chapter 8 --- Tree-Based Methods (Lecture 10).}\par\vspace{1mm}

A logistics firm records, for $n=6$ regional depots, the number of delivery vans $x$ and the
number of parcels delivered per day $y$ (in thousands):

\begin{center}
\begin{tabular}{lcccccc}
\toprule
Depot $i$ & 1 & 2 & 3 & 4 & 5 & 6\\
\midrule
Vans $x_i$   & 10 & 12 & 14 & 16 & 18 & 20\\
Parcels $y_i$ & 6  & 4  & 5  & 11 & 13 & 12\\
\bottomrule
\end{tabular}
\end{center}

A regression tree with a \emph{single} split at cutpoint $s$ partitions the depots into
$R_1=\{i: x_i<s\}$ and $R_2=\{i: x_i\ge s\}$ and predicts the region mean $\hat{y}_{R_m}$ in
each region.

\begin{enumerate}[label=(\alph*),leftmargin=7mm,itemsep=2.5mm]
  \item \pp{8} For each of the two candidate cutpoints $s=15$ and $s=17$, determine the two
    regions, their means, and the total residual sum of squares
    $\operatorname{RSS}(s)=\sum_{m=1}^{2}\sum_{i\in R_m}(y_i-\hat{y}_{R_m})^2$.
    Which cutpoint does recursive binary splitting prefer, and why?
  \item \pp{2} Suppose recursive binary splitting selects the split at $s=15$, giving regions
    $R_1=\{x<15\}$ with mean $\hat{y}_{R_1}=5$ and $R_2=\{x\ge 15\}$ with mean
    $\hat{y}_{R_2}=12$. Predict the number of parcels for a new depot with $x=13$ vans.
  \item \pp{6} A \emph{classification} tree for a different task has a terminal node containing
    $12$ training observations from three classes: $6$ ``Premium'', $3$ ``Standard'' and $3$
    ``Basic''. Compute the class proportions, the Gini index
    $G=\sum_{k}\hat{p}_{k}(1-\hat{p}_{k})$ and the entropy $D=-\sum_{k}\hat{p}_{k}\ln\hat{p}_{k}$
    of this node \emph{(use $\ln 0.5=-0.6931$ and $\ln 0.25=-1.3863$)}. Which class does the node
    predict?
  \item \pp{4} In \emph{cost-complexity pruning} one minimises, over subtrees $T\subseteq T_0$
    of the large tree $T_0$,
    \[
      \sum_{m=1}^{|T|}\ \sum_{i:\,x_i\in R_m}\bigl(y_i-\hat{y}_{R_m}\bigr)^2+\alpha\,|T| .
    \]
    Explain the role of $\alpha$: what happens for $\alpha=0$, what happens as $\alpha$ grows,
    and how is $\alpha$ chosen in practice?
\end{enumerate}

\ifdefined\withsolutions
\begin{solutionbox}
\textbf{(a) Evaluating the two cutpoints.}
For reference the root mean is $\bar{y}=(6+4+5+11+13+12)/6=51/6=8.5$.

\emph{Cutpoint $s=15$:} $R_1=\{1,2,3\}$ with $y$-values $\{6,4,5\}$, mean $\hat{y}_{R_1}=15/3=5$;
$R_2=\{4,5,6\}$ with $y$-values $\{11,13,12\}$, mean $\hat{y}_{R_2}=36/3=12$.
\[
  \begin{aligned}
  \operatorname{RSS}_{R_1}&=(6-5)^2+(4-5)^2+(5-5)^2=1+1+0=2,\\
  \operatorname{RSS}_{R_2}&=(11-12)^2+(13-12)^2+(12-12)^2=1+1+0=2,
  \end{aligned}
\]
so $\operatorname{RSS}(15)=2+2=\mathbf{4}$.

\emph{Cutpoint $s=17$:} $R_1=\{1,2,3,4\}$ with $y$-values $\{6,4,5,11\}$, mean
$\hat{y}_{R_1}=26/4=6.5$; $R_2=\{5,6\}$ with $y$-values $\{13,12\}$, mean $\hat{y}_{R_2}=12.5$.
\[
  \operatorname{RSS}_{R_1}=(6-6.5)^2+(4-6.5)^2+(5-6.5)^2+(11-6.5)^2=0.25+6.25+2.25+20.25=29,
\]
\[
  \operatorname{RSS}_{R_2}=(13-12.5)^2+(12-12.5)^2=0.25+0.25=0.5,
\]
so $\operatorname{RSS}(17)=29+0.5=\mathbf{29.5}$.

Recursive binary splitting is \emph{greedy} and takes the cutpoint with the smaller RSS: since
$4<29.5$, it prefers $\mathbf{s=15}$, which cleanly separates the three low-throughput depots
(mean $5$) from the three high-throughput depots (mean $12$).

\medskip
\textbf{(b) Prediction for $x=13$.}
$x=13<15$, so the new depot falls into $R_1$ and the tree predicts
$\hat{y}=\hat{y}_{R_1}=\mathbf{5}$ (i.e.\ $5{,}000$ parcels per day).

\medskip
\textbf{(c) Node impurity (three classes).}
$\hat{p}_{\text{Prem}}=6/12=0.5$, $\hat{p}_{\text{Std}}=3/12=0.25$,
$\hat{p}_{\text{Basic}}=3/12=0.25$.
\[
  G=0.5(1-0.5)+0.25(1-0.25)+0.25(1-0.25)=0.25+0.1875+0.1875=\mathbf{0.625}.
\]
\[
  \begin{aligned}
  D&=-\bigl(0.5\ln 0.5+0.25\ln 0.25+0.25\ln 0.25\bigr)\\
   &=-\bigl(0.5(-0.6931)+0.5(-1.3863)\bigr)=0.3466+0.6931=\mathbf{1.040}.
  \end{aligned}
\]
The node predicts the majority class, \textbf{Premium} (estimated probability $0.5$). Both $G$
and $D$ would be $0$ for a pure node; here the node is quite impure, consistent with the split
carrying three classes.

\medskip
\textbf{(d) Role of $\alpha$.}
$\alpha\ge 0$ is the price charged per terminal node, trading the subtree's fit (first term)
against its size $|T|$.
For $\alpha=0$ the criterion reduces to the training RSS, so the full unpruned tree $T_0$ is
optimal. As $\alpha$ increases, each extra leaf must ``earn its keep''; a branch whose RSS
reduction is smaller than $\alpha$ is pruned away, so the optimal subtree shrinks and one
obtains a \emph{nested sequence} of subtrees as $\alpha$ climbs (weakest-link pruning). In
practice $\alpha$ is selected by \emph{cross-validation} --- pick the $\alpha$ giving the lowest
CV error, then refit that subtree on the full data --- which balances bias and variance rather
than merely minimising training error.
\end{solutionbox}
\fi

% ============================================================
\problem{4: Ensembles --- bagging, random forests, boosting}{16}
{\small\itshape Lecture slides: Chapter 8 --- Tree-Based Methods (Lecture 10).}\par\vspace{1mm}

\begin{enumerate}[label=(\alph*),leftmargin=7mm,itemsep=2.5mm]
  \item \pp{6} Let $\hat{f}_1(x),\dots,\hat{f}_B(x)$ be $B$ identically distributed (but not
    independent) tree predictions at a fixed point $x$, each with variance $\sigma^2$ and
    pairwise correlation $\rho$. The variance of the bagged average is
    \[
      \operatorname{Var}\Bigl(\tfrac{1}{B}\sum_{b=1}^{B}\hat{f}_b(x)\Bigr)
      =\rho\,\sigma^{2}+\frac{1-\rho}{B}\,\sigma^{2}.
    \]
    (i) Explain briefly why this formula shows that bagging reduces variance and which term
    limits the reduction. (ii) Evaluate the variance for $B=500$, $\rho=0.7$ and $\sigma^2=4$,
    and compare it with the variance of a single tree. (iii) What value does the variance
    approach as $B\to\infty$?
  \item \pp{4} How do \emph{random forests} modify bagging at each split, and what is the
    recommended subset size $m$ for \emph{classification} as a function of $p$ (evaluate for
    $p=36$)? Using the variance formula
    $\operatorname{Var}=\rho\sigma^2+\tfrac{1-\rho}{B}\sigma^2$ from part~(a), explain why this
    modification improves on bagging.
  \item \pp{4} A data scientist fits a boosting model for regression trees in
    \texttt{scikit-learn}:
    \begin{center}
    \texttt{GradientBoostingRegressor(n\_estimators=500, learning\_rate=0.01, max\_depth=1)}
    \end{center}
    These three arguments are exactly the tuning parameters of boosting. State in one sentence
    each what \texttt{n\_estimators}, \texttt{learning\_rate} and \texttt{max\_depth} control,
    and explain the ``slow learning'' idea behind the small \texttt{learning\_rate}.
  \item \pp{2} What does the \emph{out-of-bag} (OOB) error of a bagged model estimate, and why
    is it essentially free to obtain?
\end{enumerate}

\ifdefined\withsolutions
\begin{solutionbox}
\textbf{(a) Variance of the bagged average.}
(i) The second term $\frac{1-\rho}{B}\sigma^2$ shrinks to $0$ as $B$ grows, so averaging over
many trees removes the ``independent part'' of their variability --- this is the variance
reduction of bagging. The first term $\rho\sigma^2$ does \emph{not} depend on $B$: the
bootstrap trees are grown on overlapping resamples of the same data (and share the same strong
predictors), so they are positively correlated, and this common part of the variance cannot be
averaged away. It is the floor that limits the reduction.
(ii) For $B=500$, $\rho=0.7$, $\sigma^2=4$:
\[
  \operatorname{Var}=0.7\cdot 4+\frac{1-0.7}{500}\cdot 4
  =2.8+\frac{0.3\cdot 4}{500}=2.8+0.0024=\mathbf{2.802},
\]
versus $\sigma^2=4$ for a single tree --- a reduction of about $30\%$.
(iii) As $B\to\infty$ the variance approaches the floor $\rho\sigma^2=0.7\cdot 4=\mathbf{2.8}$.
Enlarging $B$ further does not overfit; it only stabilises the average.

\medskip
\textbf{(b) Random forests.}
At each split the tree may consider only a fresh \emph{random subset} of $m$ out of the $p$
predictors as candidates. For classification the usual choice is $m\approx\sqrt{p}$, so for
$p=36$: $m=\sqrt{36}=\mathbf{6}$.
Why it helps: in plain bagging a dominant predictor is chosen at the top split of nearly every
tree, making the trees look alike and pushing $\rho$ up. Withholding that predictor from about
$(p-m)/p$ of the split searches \emph{decorrelates} the trees, i.e.\ lowers $\rho$; by (a) this
lowers the irreducible floor $\rho\sigma^2$ of the ensemble variance, which plain averaging can
never reach.

\medskip
\textbf{(c) Boosting tuning parameters (read from the call).}
\begin{itemize}[nosep,leftmargin=5mm]
  \item \texttt{n\_estimators}$=B=500$: the number of small trees added in sequence; unlike
    bagging, boosting \emph{can} overfit if $B$ is too large, so $B$ is tuned by cross-validation.
  \item \texttt{learning\_rate}$=\lambda=0.01$ (shrinkage): scales the contribution of each new
    tree and hence the speed of learning; a smaller $\lambda$ usually requires a larger $B$.
  \item \texttt{max\_depth}$=d=1$ (interaction depth): the number of splits per tree and the
    interaction order captured; $d=1$ means \emph{stumps} (an additive model).
\end{itemize}
\emph{Slow learning:} each new tree is fit to the current \emph{residuals}, and only the small
fraction \texttt{learning\_rate}$=\lambda$ of it is added to the model, so the fit improves
gradually and precisely where it is still weak, rather than fitting the data hard in a single
step --- which tends to generalise better.

\medskip
\textbf{(d) OOB error.}
Each bootstrap sample omits, on average, about one third of the observations
(``out-of-bag''). Predicting each observation only from the trees that did \emph{not} train on
it, and aggregating, gives the OOB error --- a valid estimate of the bagged model's
\emph{test} error. It is essentially free because it reuses a by-product of the bootstrap; no
separate validation set or cross-validation loop is needed.
\end{solutionbox}
\fi

% ============================================================
\problem{5: Neural networks by hand}{20}
{\small\itshape Lecture slides: Chapter 10 --- Deep Learning (Lecture 11).}\par\vspace{1mm}

Consider a small feed-forward network for a quantitative response with input
$x=(x_1,x_2)=(2,1)$, one hidden layer of two ReLU units ($g(z)=\max\{0,z\}$) and a linear
output unit:
\[
  A_1=g\bigl(-2+3\,x_1-1\,x_2\bigr),\qquad
  A_2=g\bigl(1-2\,x_1+1\,x_2\bigr),\qquad
  f(x)=-1+2\,A_1+3\,A_2 .
\]

\begin{enumerate}[label=(\alph*),leftmargin=7mm,itemsep=2.5mm]
  \item \pp{8} Carry out the forward pass step by step: compute both pre-activations, both
    activations $A_1,A_2$ and the output $f(x)$. Comment briefly on what the ReLU did to the
    second unit and why non-linear activations are essential.
  \item \pp{6} A classification network with $K=3$ classes returns the logits (output-layer
    pre-activations) $z=(z_1,z_2,z_3)=(2.0,\ 1.0,\ -1.0)$ for one observation whose true class
    is class~1. Compute all three softmax probabilities and the cross-entropy loss
    $-\ln\hat{p}_{\text{true}}$.
    \emph{(Use $e^{2}=7.389$, $e^{1}=2.718$, $e^{-1}=0.368$ and $\ln 10.475=2.349$.)}
  \item \pp{4} Count the parameters of a fully connected network with $p=6$ inputs, one hidden
    layer of $k=4$ units and an output layer of $K=3$ units, where every hidden and output unit
    carries a bias. Show the count separately for the two weight layers.
  \item \pp{2} A convolutional layer receives a $32\times 32$ image and applies filters of size
    $F=5$ with padding $P=0$ and stride $S=1$. Compute the spatial size of the output feature
    map from $\bigl(W-F+2P\bigr)/S+1$.
\end{enumerate}

\ifdefined\withsolutions
\begin{solutionbox}
\textbf{(a) Forward pass.}
Pre-activations:
\[
  z_1=-2+3\cdot 2-1\cdot 1=-2+6-1=3,\qquad
  z_2=1-2\cdot 2+1\cdot 1=1-4+1=-2 .
\]
ReLU activations:
\[
  A_1=g(3)=\max\{0,3\}=3,\qquad A_2=g(-2)=\max\{0,-2\}=0 .
\]
Output:
\[
  f(x)=-1+2\cdot 3+3\cdot 0=-1+6+0=\mathbf{5}.
\]
The ReLU \emph{switched the second unit off}: its negative pre-activation is clipped to $0$, so
unit~2 contributes nothing for this particular input (though it could for others). This
thresholding is exactly the source of non-linearity: with a purely linear $g$, a network of any
depth would collapse to a single linear function of $x$.

\medskip
\textbf{(b) Softmax and cross-entropy.}
Exponentials $e^{2.0}=7.389$, $e^{1.0}=2.718$, $e^{-1.0}=0.368$ sum to
$7.389+2.718+0.368=10.475$. Softmax probabilities:
\[
  \hat{p}_1=\frac{7.389}{10.475}=0.705,\qquad
  \hat{p}_2=\frac{2.718}{10.475}=0.259,\qquad
  \hat{p}_3=\frac{0.368}{10.475}=0.035,
\]
which sum to $0.999\approx 1$. Cross-entropy loss for the true class~1:
\[
  -\ln\hat{p}_1=-\ln\frac{e^{2}}{10.475}=-\bigl(2-\ln 10.475\bigr)=2.349-2=\mathbf{0.349}.
\]
(Equivalently $-\ln 0.705\approx 0.349$; a perfect prediction $\hat{p}_1=1$ would give loss 0.)

\medskip
\textbf{(c) Parameter count for $6\to 4\to 3$.}
Input $\to$ hidden: $6\cdot 4=24$ weights $+\,4$ biases $=28$.
Hidden $\to$ output: $4\cdot 3=12$ weights $+\,3$ biases $=15$.
Total: $28+15=\mathbf{43}$ parameters.

\medskip
\textbf{(d) Convolution output size.}
\[
  \frac{W-F+2P}{S}+1=\frac{32-5+2\cdot 0}{1}+1=\frac{27}{1}+1=28,
\]
so the output feature map is $\mathbf{28\times 28}$. With no padding a $5\times 5$ filter at
stride $1$ shrinks each spatial dimension by $F-1=4$.
\end{solutionbox}
\fi

% ============================================================
\problem{6: Multiple testing}{20}
{\small\itshape Lecture slides: Chapter 13 --- Multiple Testing (Lecture 12).}\par\vspace{1mm}

\begin{enumerate}[label=(\alph*),leftmargin=7mm,itemsep=2.5mm]
  \item \pp{4} A lab runs $m=10$ \emph{independent} hypothesis tests, each at level
    $\alpha=0.05$, and suppose all $10$ null hypotheses are true. Define the family-wise error
    rate (FWER) and compute it here \emph{(use $(0.95)^{10}=0.5987$)}. Interpret the result in
    one sentence.
  \item \pp{7} A team screens $m=6$ candidate biomarkers; the ordered $p$-values are
    \begin{center}
    \begin{tabular}{lcccccc}
    \toprule
    $j$ & 1 & 2 & 3 & 4 & 5 & 6\\
    \midrule
    $p_{(j)}$ & 0.002 & 0.009 & 0.012 & 0.015 & 0.040 & 0.300\\
    \bottomrule
    \end{tabular}
    \end{center}
    Control the FWER at $\alpha=0.05$ (i) with \emph{Bonferroni} and (ii) with the
    \emph{Holm step-down} procedure. Give the relevant thresholds, state exactly which
    hypotheses each method rejects, and explain why Holm can never reject fewer hypotheses than
    Bonferroni.
  \item \pp{6} Apply the \emph{Benjamini--Hochberg} (BH) procedure at FDR level $q=0.05$ to the
    same six ordered $p$-values ($m=6$; $p_{(1)},\dots,p_{(6)}=0.002,\ 0.009,\ 0.012,\ 0.015,\
    0.040,\ 0.300$): give the thresholds $jq/m$, determine which hypotheses are rejected, and
    state precisely what ``FDR control at $q=0.05$'' guarantees for the resulting set of
    rejections.
  \item \pp{3} A pharmaceutical company runs a \emph{confirmatory} Phase-III trial with $m=4$
    pre-specified primary endpoints; a single significant endpoint could support regulatory
    approval of the drug. Should the company control the FWER or the FDR? Justify your
    recommendation.
\end{enumerate}

\ifdefined\withsolutions
\begin{solutionbox}
\textbf{(a) FWER for $m=10$ independent tests.}
$\text{FWER}=\Pr(\text{at least one Type I error})=\Pr(V\ge 1)$, where $V$ counts the true
nulls that are falsely rejected. With independence and all nulls true,
\[
  \text{FWER}=1-\Pr(\text{no false rejection})=1-(1-0.05)^{10}=1-(0.95)^{10}
  =1-0.5987=\mathbf{0.4013}.
\]
Interpretation: although each individual test uses the $5\%$ level, there is about a $40\%$
chance of at least one false discovery across the ten tests --- testing many hypotheses without
adjustment makes spurious findings very likely.

\medskip
\textbf{(b) Bonferroni and Holm at $\alpha=0.05$.}
\emph{Bonferroni:} reject $H_{(j)}$ iff $p_{(j)}\le\alpha/m=0.05/6=0.00833$.
Only $p_{(1)}=0.002\le 0.00833$ \checkmark; already $p_{(2)}=0.009>0.00833$ $\times$.
$\Rightarrow$ Bonferroni rejects $H_{(1)}$ alone: \textbf{1 rejection}. (Note $p_{(2)}=0.009$ is
below the naive $0.05$ cutoff, yet Bonferroni does not reject it.)

\emph{Holm:} compare $p_{(j)}$ with $\alpha/(m-j+1)$ in order and stop at the first failure:
\begin{center}
\begin{tabular}{lcccccc}
\toprule
$j$ & 1 & 2 & 3 & 4 & 5 & 6\\
\midrule
$p_{(j)}$ & 0.002 & 0.009 & 0.012 & 0.015 & 0.040 & 0.300\\
$\alpha/(m-j+1)$ & 0.00833 & 0.0100 & 0.0125 & 0.0167 & 0.0250 & 0.0500\\
reject? & \checkmark & \checkmark & \checkmark & \checkmark & stop & ---\\
\bottomrule
\end{tabular}
\end{center}
$0.002\le 0.00833$, $0.009\le 0.0100$, $0.012\le 0.0125$, $0.015\le 0.0167$, but
$0.040>0.0250$: stop. $\Rightarrow$ Holm rejects $H_{(1)},H_{(2)},H_{(3)},H_{(4)}$:
\textbf{4 rejections}.

\emph{Why Holm $\supseteq$ Bonferroni:} Holm's thresholds $\alpha/(m-j+1)$ are all $\ge\alpha/m$
(with equality only at $j=1$), so any $p$-value that clears the strict Bonferroni bar also
clears Holm's more lenient later hurdles --- Holm is uniformly at least as powerful, while still
controlling the FWER at $\alpha$ with no independence assumption.

\medskip
\textbf{(c) Benjamini--Hochberg at $q=0.05$.}
Thresholds $jq/m=j\cdot 0.05/6$:
\begin{center}
\begin{tabular}{lcccccc}
\toprule
$j$ & 1 & 2 & 3 & 4 & 5 & 6\\
\midrule
$p_{(j)}$ & 0.002 & 0.009 & 0.012 & 0.015 & 0.040 & 0.300\\
$jq/m$ & 0.00833 & 0.0167 & 0.0250 & 0.0333 & 0.0417 & 0.0500\\
$p_{(j)}\le jq/m$? & yes & yes & yes & yes & \textbf{yes} & no\\
\bottomrule
\end{tabular}
\end{center}
$L=$ largest $j$ with $p_{(j)}\le jq/m$: here $j=5$ (since $0.040\le 0.0417$, while
$0.300>0.0500$). BH rejects \emph{all} $H_{(j)}$ with $j\le L$, i.e.\
$H_{(1)},\dots,H_{(5)}$: \textbf{5 rejections}.
Guarantee: $\text{FDR}=E\bigl[V/\max(R,1)\bigr]\le q=0.05$ --- \emph{on average} at most $5\%$
of the rejected hypotheses are false discoveries. It does \emph{not} promise that at most $5\%$
of these particular $5$ rejections are false, nor does it bound the probability of making any
false rejection (that would be FWER control).

\medskip
\textbf{(d) Confirmatory trial recommendation.}
Control the \textbf{FWER}. Here even a \emph{single} false positive is very costly --- it could
lead to approving an ineffective drug --- so the study must guard against any Type I error, not
merely keep their average proportion small. With only $m=4$ pre-specified endpoints, an FWER
procedure (e.g.\ Holm or Bonferroni at $0.05/4=0.0125$) sacrifices little power. FDR control,
appropriate for large exploratory screens where a few false leads are tolerable, is too
permissive for a confirmatory, decision-driving trial.
\end{solutionbox}
\fi

% ============================================================
\problem{7: Cumulative essentials (Chapters 2--6)}{20}
{\small\itshape Lecture slides: Chapters 2--6 (Lectures 1--8).}\par\vspace{1mm}

\begin{enumerate}[label=(\alph*),leftmargin=7mm,itemsep=2.5mm]
  \item \pp{4} A medical screening classifier is evaluated on $200$ patients (label
    \texttt{1}$=$diseased, \texttt{0}$=$healthy). In \texttt{scikit-learn},
    \texttt{confusion\_matrix(y\_true, y\_pred)} returns the array below (rows $=$ true class,
    columns $=$ predicted class; label order $[0,1]$):
    \begin{center}
    \begin{tabular}{lcc}
    \toprule
     & predicted \texttt{0} & predicted \texttt{1}\\
    \midrule
    true \texttt{0} (healthy)  & $120$ & $30$\\
    true \texttt{1} (diseased) & $10$  & $40$\\
    \bottomrule
    \end{tabular}
    \end{center}
    Read off the counts and compute the overall error rate, the sensitivity, the specificity and
    the precision (positive predictive value) for the positive class \texttt{1}. Which single
    number would trace out the ROC curve if it were varied, and in which direction would
    sensitivity and specificity move?
  \item \pp{4} Consider $K$-nearest-neighbours. (i) Using the training set
    $(x,y)$: $(1,3)$, $(4,6)$, $(6,8)$, $(7,7)$, $(10,12)$, predict the response at $x_0=5$ with $k=3$
    (KNN regression, unweighted average). (ii) Explain how $k$ governs the bias--variance
    trade-off, which extreme ($k=1$ or $k$ large) is the most flexible, and why the
    \emph{training} error of $1$-NN is always $0$.
  \item \pp{4} Explain in one or two sentences what $k$-fold cross-validation does, then give
    one reason to prefer $k=5$ or $10$ over \emph{LOOCV} and one reason to prefer it over a
    \emph{single} validation split.
  \item \pp{2} You have $p=2{,}000$ predictors but expect only a handful to be truly related to
    the response. Ridge or lasso --- which penalty is more appropriate, and where would
    principal-components regression (PCR) sit relative to your choice?
  \item \pp{6} Match each scenario to exactly one method from
    \{\emph{linear regression}, \emph{lasso}, \emph{logistic regression},
    \emph{random forest}\} (each method used exactly once) and justify each choice in one
    sentence:
    \begin{enumerate}[label=(\roman*),nosep,topsep=1mm]
      \item Estimate how much one additional year of schooling adds to annual earnings, with a
        confidence interval, from $n=5{,}000$ survey responses.
      \item Predict a continuous blood-pressure outcome from $p=8{,}000$ genetic markers for
        $n=150$ patients, and report a short list of the relevant markers.
      \item Classify each incoming e-mail as spam or not, and explain to auditors how each
        feature changes the odds of an e-mail being spam.
      \item Predict machine failure from hundreds of mixed sensor variables with strong
        non-linear effects and interactions; only predictive accuracy matters.
    \end{enumerate}
\end{enumerate}

\ifdefined\withsolutions
\begin{solutionbox}
\textbf{(a) Reading the \texttt{confusion\_matrix} output.}
The array is ordered $[0,1]$ with rows $=$ true and columns $=$ predicted, so the diagonal holds
the correct predictions: from the table $TN=120$ and $FP=30$ (true \texttt{0}), $FN=10$ and
$TP=40$ (true \texttt{1}); total $200$, with $50$ diseased and $150$ healthy.
\[
  \text{error rate}=\frac{FP+FN}{200}=\frac{30+10}{200}=\mathbf{0.200},\qquad
  \text{sensitivity}=\frac{TP}{TP+FN}=\frac{40}{50}=\mathbf{0.800},
\]
\[
  \text{specificity}=\frac{TN}{TN+FP}=\frac{120}{150}=\mathbf{0.800},\qquad
  \text{precision}=\frac{TP}{TP+FP}=\frac{40}{70}=\mathbf{0.571}.
\]
Varying the \emph{classification threshold} on the predicted probability traces out the ROC
curve. Lowering the threshold labels more patients ``diseased'', which \emph{raises}
sensitivity (fewer missed cases) but \emph{lowers} specificity (more false alarms); raising it
does the reverse --- the ROC curve captures exactly this trade-off.

\medskip
\textbf{(b) $K$-nearest-neighbours.}
(i) Distances from $x_0=5$ are $|1-5|=4$, $|4-5|=1$, $|6-5|=1$, $|7-5|=2$, $|10-5|=5$; the three
nearest neighbours are $x=4,6,7$ with responses $6,8,7$, so
$\hat{y}(5)=(6+8+7)/3=21/3=\mathbf{7}$.
(ii) Small $k$ gives a very flexible fit: low bias but high variance, since the prediction hangs
on a few nearby points and a wiggly boundary. Large $k$ averages over many neighbours, giving a
smooth, stable fit: high bias but low variance. The most flexible extreme is $\mathbf{k=1}$
(flexibility $\propto 1/k$). The training error of $1$-NN is always $0$ because each training
point's single nearest neighbour is \emph{itself}, so it predicts its own observed response
exactly --- which is why training error is no guide to test performance.

\medskip
\textbf{(c) $k$-fold cross-validation.}
It splits the data into $k$ roughly equal folds; each fold in turn serves as the validation set
while the model is trained on the other $k-1$ folds, and the $k$ resulting errors are averaged.
\emph{vs.\ LOOCV:} the $n$ LOOCV fits use almost identical training sets, so their errors are
highly correlated and their average has \emph{higher variance} than $5$- or $10$-fold CV (and
LOOCV costs $n$ fits, absent shortcut formulas).
\emph{vs.\ a single validation split:} a single split gives a high-variance estimate that
depends heavily on the particular partition and is biased upward because only part of the data
trains the model; $k$-fold CV averages over $k$ folds and trains on $(k-1)/k$ of the data,
giving a more stable, less biased estimate.

\medskip
\textbf{(d) Lasso.}
The $\ell_1$ penalty drives many coefficients \emph{exactly to zero} and thus performs variable
selection --- ideal when the truth is sparse (a few real signals among $2{,}000$ predictors) and
an interpretable small model is wanted. Ridge only shrinks coefficients and keeps all $2{,}000$
predictors, so it tends to lose here. PCR is likewise \emph{not} sparse: it regresses on a few
principal components, so all original variables load on the retained components and no variable
is truly dropped --- like ridge it does not perform selection, so it does not deliver the short
interpretable list that lasso does.

\medskip
\textbf{(e) Matching.}
\begin{itemize}[nosep,leftmargin=5mm]
  \item (i) \textbf{Linear regression}: a quantitative response with $n\gg p$ and the goal of
    \emph{inference} --- an effect size with a confidence interval --- exactly what OLS with
    standard errors provides.
  \item (ii) \textbf{Lasso}: $p\gg n$ makes OLS infeasible; the $\ell_1$ penalty regularises the
    fit and returns a short list of relevant markers.
  \item (iii) \textbf{Logistic regression}: a binary response with interpretable coefficients
    ($e^{\hat\beta_j}$ are odds ratios) --- well suited to explaining odds to auditors.
  \item (iv) \textbf{Random forest}: automatically captures non-linearities and interactions,
    is accurate and robust with many mixed predictors, and interpretability is not required.
\end{itemize}
\end{solutionbox}
\fi

\vspace{6mm}
\begin{center}
\rule{0.35\textwidth}{0.4pt}\\[1mm]
\textit{End of the examination --- good luck!}
\end{center}

\end{document}
"
% ============================================================
\documentclass[11pt,a4paper]{article}
\usepackage[utf8]{inputenc}
\usepackage[T1]{fontenc}
\usepackage[english]{babel}
\usepackage[margin=2.2cm]{geometry}
\usepackage{amsmath,amssymb}
\usepackage{booktabs,array}
\usepackage{enumitem}
\usepackage{xcolor}
\usepackage{tcolorbox}
\tcbuselibrary{skins,breakable}

\setlength{\parindent}{0pt}

% Teal solution box (only shown when \withsolutions is defined)
\newtcolorbox{solutionbox}{enhanced,breakable,colback=teal!5,colframe=teal!55!black,
  boxrule=0pt,leftrule=2.5pt,arc=2pt,fonttitle=\bfseries,title=Solution,
  left=2mm,right=2mm,top=1mm,bottom=1.5mm}

\newcommand{\problem}[2]{\par\vspace{7mm}\noindent{\large\textbf{Problem #1}}%
  \hfill{\large\textbf{(#2 points)}}\par\nopagebreak\vspace{0.6mm}%
  \noindent\rule{\textwidth}{0.6pt}\par\nopagebreak\vspace{2.5mm}}
\newcommand{\pp}[1]{\textbf{[#1~P]}\enspace}

\begin{document}

% ------------------------------------------------------------
% Header block
% ------------------------------------------------------------
\begin{center}
  {\Large\textbf{HSBI --- Bielefeld University of Applied Sciences and Arts}}\\[1.2mm]
  {\large Quantitative Research Methods}\\[1mm]
  {\normalsize Summer Semester 2026 \quad$\cdot$\quad Examiner: Prof.\ Dr.\ Christoph Weisser}\\[4.5mm]
  {\LARGE\textbf{Final Mock Examination --- Set B (Lectures 1--12)}}\\[1.2mm]
  {\normalsize Comprehensive coverage of ISLP Chapters 2--8, 10 and 13,\\
   with emphasis on Chapters 7 (splines/GAMs), 8 (trees), 10 (deep learning) and 13 (multiple testing)}%
  \ifdefined\withsolutions\\[2.2mm]{\large\textcolor{teal!55!black}{\textbf{--- Version with detailed solutions ---}}}\fi
\end{center}

\vspace{2.5mm}
\begin{center}
\begin{tabular}{@{}ll@{}}
  \textbf{Duration:} & \textbf{120 minutes}\\
  \textbf{Total credit:} & \textbf{120 points}\\
  \textbf{Allowed aids:} & non-programmable calculator; one handwritten A4 sheet of notes (both sides)\\
\end{tabular}
\end{center}

\vspace{3mm}
Name: \rule[-1mm]{0.38\textwidth}{0.4pt}\hfill Matriculation no.: \rule[-1mm]{0.24\textwidth}{0.4pt}

\vspace{4.5mm}
\begin{center}
\renewcommand{\arraystretch}{1.35}
\begin{tabular}{|l|c|c|c|c|c|c|c|c|}
\hline
\textbf{Problem} & \textbf{P1} & \textbf{P2} & \textbf{P3} & \textbf{P4} & \textbf{P5} & \textbf{P6} & \textbf{P7} & \textbf{Total}\\\hline
Maximum points & 16 & 8 & 20 & 16 & 20 & 20 & 20 & \textbf{120}\\\hline
Achieved points & \hspace{8mm} & \hspace{8mm} & \hspace{8mm} & \hspace{8mm} & \hspace{8mm} & \hspace{8mm} & \hspace{8mm} & \hspace{10mm}\\\hline
\end{tabular}
\end{center}

\vspace{2.5mm}
\textbf{Instructions}
\begin{itemize}[nosep,leftmargin=5.5mm]
  \item Justify all answers. Numerical results without a derivation or a brief reasoning receive no credit.
  \item Round final numerical results to 3 decimal places unless stated otherwise. Small deviations caused by carrying rounded intermediate results are accepted.
  \item All numerical constants you need (powers of $e$, logarithms, $(0.95)^{10}$, $t$-quantiles) are provided in the problems.
  \item The number of points per problem roughly equals the recommended working time in minutes.
\end{itemize}
\vspace{1mm}\noindent\rule{\textwidth}{0.6pt}

% ============================================================
\problem{1: Moving beyond linearity --- splines}{16}
{\small\itshape Lecture slides: Chapter 7 --- Moving Beyond Linearity (Lecture 9).}\par\vspace{1mm}

Throughout this problem, consider regression splines built on a single predictor $X$ using
$K=5$ interior knots.

\begin{enumerate}[label=(\alph*),leftmargin=7mm,itemsep=2.5mm]
  \item \pp{4} A \emph{cubic spline} is a piecewise cubic polynomial that is continuous and has
    continuous first and second derivatives at every knot. How many degrees of freedom
    (free parameters, counting the intercept) does a cubic spline with $K$ interior knots use?
    Justify your answer by counting the parameters of the polynomial pieces and subtracting the
    constraints at the knots, and evaluate for $K=5$.
  \item \pp{4} How many degrees of freedom does a \emph{natural} cubic spline with the same
    $K=5$ knots use? State the additional constraint a natural spline imposes, in which region
    of the predictor space it acts, and why practitioners find this constraint useful.
  \item \pp{4} A \emph{smoothing spline} is the function $g$ that minimises
    \[
      \sum_{i=1}^{n}\bigl(y_i-g(x_i)\bigr)^2+\lambda\int g''(t)^2\,dt .
    \]
    Explain the meaning of the two terms, and describe the shape of the solution $\hat{g}$ in
    the two limiting cases $\lambda\to 0$ and $\lambda\to\infty$, giving the effective degrees
    of freedom that $\hat{g}$ approaches in each case.
  \item \pp{4} Write down the \emph{truncated-power basis} representation of a cubic spline with
    a single knot at $\xi$, define the truncated-power basis function $(x-\xi)^3_+$ explicitly,
    and explain in one sentence why adding this term keeps the fit continuous with continuous
    first and second derivatives at $\xi$.
\end{enumerate}

\ifdefined\withsolutions
\begin{solutionbox}
\textbf{(a) Degrees of freedom of a cubic spline.}
With $K$ interior knots the range of $X$ is divided into $K+1$ intervals, and on each interval
the fit is a cubic with $4$ coefficients, giving $4(K+1)=4K+4$ free parameters. At every knot
the requirement that $g$, $g'$ and $g''$ join continuously imposes $3$ constraints, i.e.\ $3K$
constraints in total. Hence
\[
  \mathrm{df}=4(K+1)-3K=K+4 .
\]
For $K=5$: $\mathrm{df}=5+4=\mathbf{9}$ (equivalently the basis intercept, $x$, $x^2$, $x^3$
plus one truncated-power term per knot, $4+K$).

\medskip
\textbf{(b) Natural cubic spline.}
A natural spline adds the requirement that the function be \emph{linear beyond the two boundary
knots}, i.e.\ on the outer regions $X<\xi_1$ and $X>\xi_K$. Forcing the quadratic and cubic
terms to vanish in each of those two regions removes $2$ parameters per side, $4$ in total:
\[
  \mathrm{df}=(K+4)-4=K,\qquad K=5:\ \mathrm{df}=\mathbf{5}.
\]
Usefulness: ordinary polynomials and unconstrained splines behave erratically near the
boundaries, where the data are sparse; forcing linearity there reduces the \emph{variance} of
predictions at and beyond the outer knots, at the cost of only a little extra bias.

\medskip
\textbf{(c) Smoothing spline.}
The first term (the RSS) rewards closeness to the observed data, while the second term
penalises \emph{roughness}: $g''$ measures the local curvature of $g$, so $\int g''(t)^2dt$ is
large for wiggly functions, and $\lambda\ge 0$ controls the trade-off.
\emph{$\lambda\to 0$:} the penalty vanishes, so $\hat{g}$ may be arbitrarily wiggly and
\emph{interpolates} every training point (zero training RSS, extreme overfitting); the
effective degrees of freedom approach $n$.
\emph{$\lambda\to\infty$:} any curvature becomes infinitely costly, forcing $g''\equiv 0$, so
$\hat{g}$ is a straight line, namely the ordinary \emph{least-squares regression line}
(effective df $\to 2$). For intermediate $\lambda$ (chosen e.g.\ by LOOCV) the effective df
decrease smoothly from $n$ down to $2$.

\medskip
\textbf{(d) Truncated-power basis, single knot $\xi$.}
\[
  f(x)=\beta_0+\beta_1x+\beta_2x^2+\beta_3x^3+\beta_4\,(x-\xi)^3_+,
  \qquad
  (x-\xi)^3_+=\begin{cases}(x-\xi)^3, & x>\xi\\[0.5mm] 0, & x\le\xi\end{cases}
\]
At $x=\xi$ the added function $(x-\xi)^3_+$ and its first and second derivatives are all $0$, so
including it alters only the \emph{third} derivative across the knot; the fit therefore stays
continuous with continuous first and second derivatives --- precisely the smoothness that
defines a cubic spline.
\end{solutionbox}
\fi

% ============================================================
\problem{2: Generalised additive models}{8}
{\small\itshape Lecture slides: Chapter 7 --- Moving Beyond Linearity (Lecture 9).}\par\vspace{1mm}

\begin{enumerate}[label=(\alph*),leftmargin=7mm,itemsep=2.5mm]
  \item \pp{3} Write down the general form of a \emph{generalised additive model} (GAM) for a
    quantitative response $Y$ and predictors $X_1,\dots,X_p$, and explain in one sentence how it
    generalises multiple linear regression.
  \item \pp{3} State one important \emph{interpretability advantage} of GAMs (concerning the
    fitted functions $f_j$) and their central structural \emph{limitation}, including how that
    limitation can be repaired.
  \item \pp{2} Describe in one or two sentences how the \emph{backfitting} algorithm fits a GAM.
\end{enumerate}

\ifdefined\withsolutions
\begin{solutionbox}
\textbf{(a) Model form.}
\[
  Y=\beta_0+f_1(X_1)+f_2(X_2)+\dots+f_p(X_p)+\varepsilon .
\]
Multiple linear regression is the special case in which every $f_j(X_j)=\beta_jX_j$; a GAM
replaces each straight-line term by a flexible smooth function $f_j$ (e.g.\ a natural or
smoothing spline, or a local regression) while retaining the \emph{additive} structure.

\medskip
\textbf{(b) Advantage and limitation.}
\emph{Advantage (interpretability):} because the effects add up, the contribution of each
predictor is carried by its own function $f_j$, which can be plotted and read off one panel at
a time with the other predictors held fixed --- non-linear effects stay as easy to inspect as
coefficients in a linear model.
\emph{Limitation:} pure additivity rules out \emph{interactions} --- the effect of $X_j$ is not
allowed to depend on the value of another $X_k$. Repair: introduce interaction terms by hand,
e.g.\ a product $X_j\times X_k$ or a low-dimensional smooth surface $f_{jk}(X_j,X_k)$ as an
extra additive component.

\medskip
\textbf{(c) Backfitting.}
Loop over the predictors and repeatedly re-estimate one $f_j$ at a time (with a smoother) using
the \emph{partial residuals} $r_i=y_i-\beta_0-\sum_{k\ne j}f_k(x_{ik})$, keeping the other
fitted functions fixed; repeat the cycles until the $f_j$ no longer change appreciably.
\end{solutionbox}
\fi

% ============================================================
\problem{3: Regression and classification trees by hand}{20}
{\small\itshape Lecture slides: Chapter 8 --- Tree-Based Methods (Lecture 10).}\par\vspace{1mm}

A logistics firm records, for $n=6$ regional depots, the number of delivery vans $x$ and the
number of parcels delivered per day $y$ (in thousands):

\begin{center}
\begin{tabular}{lcccccc}
\toprule
Depot $i$ & 1 & 2 & 3 & 4 & 5 & 6\\
\midrule
Vans $x_i$   & 10 & 12 & 14 & 16 & 18 & 20\\
Parcels $y_i$ & 6  & 4  & 5  & 11 & 13 & 12\\
\bottomrule
\end{tabular}
\end{center}

A regression tree with a \emph{single} split at cutpoint $s$ partitions the depots into
$R_1=\{i: x_i<s\}$ and $R_2=\{i: x_i\ge s\}$ and predicts the region mean $\hat{y}_{R_m}$ in
each region.

\begin{enumerate}[label=(\alph*),leftmargin=7mm,itemsep=2.5mm]
  \item \pp{8} For each of the two candidate cutpoints $s=15$ and $s=17$, determine the two
    regions, their means, and the total residual sum of squares
    $\operatorname{RSS}(s)=\sum_{m=1}^{2}\sum_{i\in R_m}(y_i-\hat{y}_{R_m})^2$.
    Which cutpoint does recursive binary splitting prefer, and why?
  \item \pp{2} Suppose recursive binary splitting selects the split at $s=15$, giving regions
    $R_1=\{x<15\}$ with mean $\hat{y}_{R_1}=5$ and $R_2=\{x\ge 15\}$ with mean
    $\hat{y}_{R_2}=12$. Predict the number of parcels for a new depot with $x=13$ vans.
  \item \pp{6} A \emph{classification} tree for a different task has a terminal node containing
    $12$ training observations from three classes: $6$ ``Premium'', $3$ ``Standard'' and $3$
    ``Basic''. Compute the class proportions, the Gini index
    $G=\sum_{k}\hat{p}_{k}(1-\hat{p}_{k})$ and the entropy $D=-\sum_{k}\hat{p}_{k}\ln\hat{p}_{k}$
    of this node \emph{(use $\ln 0.5=-0.6931$ and $\ln 0.25=-1.3863$)}. Which class does the node
    predict?
  \item \pp{4} In \emph{cost-complexity pruning} one minimises, over subtrees $T\subseteq T_0$
    of the large tree $T_0$,
    \[
      \sum_{m=1}^{|T|}\ \sum_{i:\,x_i\in R_m}\bigl(y_i-\hat{y}_{R_m}\bigr)^2+\alpha\,|T| .
    \]
    Explain the role of $\alpha$: what happens for $\alpha=0$, what happens as $\alpha$ grows,
    and how is $\alpha$ chosen in practice?
\end{enumerate}

\ifdefined\withsolutions
\begin{solutionbox}
\textbf{(a) Evaluating the two cutpoints.}
For reference the root mean is $\bar{y}=(6+4+5+11+13+12)/6=51/6=8.5$.

\emph{Cutpoint $s=15$:} $R_1=\{1,2,3\}$ with $y$-values $\{6,4,5\}$, mean $\hat{y}_{R_1}=15/3=5$;
$R_2=\{4,5,6\}$ with $y$-values $\{11,13,12\}$, mean $\hat{y}_{R_2}=36/3=12$.
\[
  \begin{aligned}
  \operatorname{RSS}_{R_1}&=(6-5)^2+(4-5)^2+(5-5)^2=1+1+0=2,\\
  \operatorname{RSS}_{R_2}&=(11-12)^2+(13-12)^2+(12-12)^2=1+1+0=2,
  \end{aligned}
\]
so $\operatorname{RSS}(15)=2+2=\mathbf{4}$.

\emph{Cutpoint $s=17$:} $R_1=\{1,2,3,4\}$ with $y$-values $\{6,4,5,11\}$, mean
$\hat{y}_{R_1}=26/4=6.5$; $R_2=\{5,6\}$ with $y$-values $\{13,12\}$, mean $\hat{y}_{R_2}=12.5$.
\[
  \operatorname{RSS}_{R_1}=(6-6.5)^2+(4-6.5)^2+(5-6.5)^2+(11-6.5)^2=0.25+6.25+2.25+20.25=29,
\]
\[
  \operatorname{RSS}_{R_2}=(13-12.5)^2+(12-12.5)^2=0.25+0.25=0.5,
\]
so $\operatorname{RSS}(17)=29+0.5=\mathbf{29.5}$.

Recursive binary splitting is \emph{greedy} and takes the cutpoint with the smaller RSS: since
$4<29.5$, it prefers $\mathbf{s=15}$, which cleanly separates the three low-throughput depots
(mean $5$) from the three high-throughput depots (mean $12$).

\medskip
\textbf{(b) Prediction for $x=13$.}
$x=13<15$, so the new depot falls into $R_1$ and the tree predicts
$\hat{y}=\hat{y}_{R_1}=\mathbf{5}$ (i.e.\ $5{,}000$ parcels per day).

\medskip
\textbf{(c) Node impurity (three classes).}
$\hat{p}_{\text{Prem}}=6/12=0.5$, $\hat{p}_{\text{Std}}=3/12=0.25$,
$\hat{p}_{\text{Basic}}=3/12=0.25$.
\[
  G=0.5(1-0.5)+0.25(1-0.25)+0.25(1-0.25)=0.25+0.1875+0.1875=\mathbf{0.625}.
\]
\[
  \begin{aligned}
  D&=-\bigl(0.5\ln 0.5+0.25\ln 0.25+0.25\ln 0.25\bigr)\\
   &=-\bigl(0.5(-0.6931)+0.5(-1.3863)\bigr)=0.3466+0.6931=\mathbf{1.040}.
  \end{aligned}
\]
The node predicts the majority class, \textbf{Premium} (estimated probability $0.5$). Both $G$
and $D$ would be $0$ for a pure node; here the node is quite impure, consistent with the split
carrying three classes.

\medskip
\textbf{(d) Role of $\alpha$.}
$\alpha\ge 0$ is the price charged per terminal node, trading the subtree's fit (first term)
against its size $|T|$.
For $\alpha=0$ the criterion reduces to the training RSS, so the full unpruned tree $T_0$ is
optimal. As $\alpha$ increases, each extra leaf must ``earn its keep''; a branch whose RSS
reduction is smaller than $\alpha$ is pruned away, so the optimal subtree shrinks and one
obtains a \emph{nested sequence} of subtrees as $\alpha$ climbs (weakest-link pruning). In
practice $\alpha$ is selected by \emph{cross-validation} --- pick the $\alpha$ giving the lowest
CV error, then refit that subtree on the full data --- which balances bias and variance rather
than merely minimising training error.
\end{solutionbox}
\fi

% ============================================================
\problem{4: Ensembles --- bagging, random forests, boosting}{16}
{\small\itshape Lecture slides: Chapter 8 --- Tree-Based Methods (Lecture 10).}\par\vspace{1mm}

\begin{enumerate}[label=(\alph*),leftmargin=7mm,itemsep=2.5mm]
  \item \pp{6} Let $\hat{f}_1(x),\dots,\hat{f}_B(x)$ be $B$ identically distributed (but not
    independent) tree predictions at a fixed point $x$, each with variance $\sigma^2$ and
    pairwise correlation $\rho$. The variance of the bagged average is
    \[
      \operatorname{Var}\Bigl(\tfrac{1}{B}\sum_{b=1}^{B}\hat{f}_b(x)\Bigr)
      =\rho\,\sigma^{2}+\frac{1-\rho}{B}\,\sigma^{2}.
    \]
    (i) Explain briefly why this formula shows that bagging reduces variance and which term
    limits the reduction. (ii) Evaluate the variance for $B=500$, $\rho=0.7$ and $\sigma^2=4$,
    and compare it with the variance of a single tree. (iii) What value does the variance
    approach as $B\to\infty$?
  \item \pp{4} How do \emph{random forests} modify bagging at each split, and what is the
    recommended subset size $m$ for \emph{classification} as a function of $p$ (evaluate for
    $p=36$)? Using the variance formula
    $\operatorname{Var}=\rho\sigma^2+\tfrac{1-\rho}{B}\sigma^2$ from part~(a), explain why this
    modification improves on bagging.
  \item \pp{4} A data scientist fits a boosting model for regression trees in
    \texttt{scikit-learn}:
    \begin{center}
    \texttt{GradientBoostingRegressor(n\_estimators=500, learning\_rate=0.01, max\_depth=1)}
    \end{center}
    These three arguments are exactly the tuning parameters of boosting. State in one sentence
    each what \texttt{n\_estimators}, \texttt{learning\_rate} and \texttt{max\_depth} control,
    and explain the ``slow learning'' idea behind the small \texttt{learning\_rate}.
  \item \pp{2} What does the \emph{out-of-bag} (OOB) error of a bagged model estimate, and why
    is it essentially free to obtain?
\end{enumerate}

\ifdefined\withsolutions
\begin{solutionbox}
\textbf{(a) Variance of the bagged average.}
(i) The second term $\frac{1-\rho}{B}\sigma^2$ shrinks to $0$ as $B$ grows, so averaging over
many trees removes the ``independent part'' of their variability --- this is the variance
reduction of bagging. The first term $\rho\sigma^2$ does \emph{not} depend on $B$: the
bootstrap trees are grown on overlapping resamples of the same data (and share the same strong
predictors), so they are positively correlated, and this common part of the variance cannot be
averaged away. It is the floor that limits the reduction.
(ii) For $B=500$, $\rho=0.7$, $\sigma^2=4$:
\[
  \operatorname{Var}=0.7\cdot 4+\frac{1-0.7}{500}\cdot 4
  =2.8+\frac{0.3\cdot 4}{500}=2.8+0.0024=\mathbf{2.802},
\]
versus $\sigma^2=4$ for a single tree --- a reduction of about $30\%$.
(iii) As $B\to\infty$ the variance approaches the floor $\rho\sigma^2=0.7\cdot 4=\mathbf{2.8}$.
Enlarging $B$ further does not overfit; it only stabilises the average.

\medskip
\textbf{(b) Random forests.}
At each split the tree may consider only a fresh \emph{random subset} of $m$ out of the $p$
predictors as candidates. For classification the usual choice is $m\approx\sqrt{p}$, so for
$p=36$: $m=\sqrt{36}=\mathbf{6}$.
Why it helps: in plain bagging a dominant predictor is chosen at the top split of nearly every
tree, making the trees look alike and pushing $\rho$ up. Withholding that predictor from about
$(p-m)/p$ of the split searches \emph{decorrelates} the trees, i.e.\ lowers $\rho$; by (a) this
lowers the irreducible floor $\rho\sigma^2$ of the ensemble variance, which plain averaging can
never reach.

\medskip
\textbf{(c) Boosting tuning parameters (read from the call).}
\begin{itemize}[nosep,leftmargin=5mm]
  \item \texttt{n\_estimators}$=B=500$: the number of small trees added in sequence; unlike
    bagging, boosting \emph{can} overfit if $B$ is too large, so $B$ is tuned by cross-validation.
  \item \texttt{learning\_rate}$=\lambda=0.01$ (shrinkage): scales the contribution of each new
    tree and hence the speed of learning; a smaller $\lambda$ usually requires a larger $B$.
  \item \texttt{max\_depth}$=d=1$ (interaction depth): the number of splits per tree and the
    interaction order captured; $d=1$ means \emph{stumps} (an additive model).
\end{itemize}
\emph{Slow learning:} each new tree is fit to the current \emph{residuals}, and only the small
fraction \texttt{learning\_rate}$=\lambda$ of it is added to the model, so the fit improves
gradually and precisely where it is still weak, rather than fitting the data hard in a single
step --- which tends to generalise better.

\medskip
\textbf{(d) OOB error.}
Each bootstrap sample omits, on average, about one third of the observations
(``out-of-bag''). Predicting each observation only from the trees that did \emph{not} train on
it, and aggregating, gives the OOB error --- a valid estimate of the bagged model's
\emph{test} error. It is essentially free because it reuses a by-product of the bootstrap; no
separate validation set or cross-validation loop is needed.
\end{solutionbox}
\fi

% ============================================================
\problem{5: Neural networks by hand}{20}
{\small\itshape Lecture slides: Chapter 10 --- Deep Learning (Lecture 11).}\par\vspace{1mm}

Consider a small feed-forward network for a quantitative response with input
$x=(x_1,x_2)=(2,1)$, one hidden layer of two ReLU units ($g(z)=\max\{0,z\}$) and a linear
output unit:
\[
  A_1=g\bigl(-2+3\,x_1-1\,x_2\bigr),\qquad
  A_2=g\bigl(1-2\,x_1+1\,x_2\bigr),\qquad
  f(x)=-1+2\,A_1+3\,A_2 .
\]

\begin{enumerate}[label=(\alph*),leftmargin=7mm,itemsep=2.5mm]
  \item \pp{8} Carry out the forward pass step by step: compute both pre-activations, both
    activations $A_1,A_2$ and the output $f(x)$. Comment briefly on what the ReLU did to the
    second unit and why non-linear activations are essential.
  \item \pp{6} A classification network with $K=3$ classes returns the logits (output-layer
    pre-activations) $z=(z_1,z_2,z_3)=(2.0,\ 1.0,\ -1.0)$ for one observation whose true class
    is class~1. Compute all three softmax probabilities and the cross-entropy loss
    $-\ln\hat{p}_{\text{true}}$.
    \emph{(Use $e^{2}=7.389$, $e^{1}=2.718$, $e^{-1}=0.368$ and $\ln 10.475=2.349$.)}
  \item \pp{4} Count the parameters of a fully connected network with $p=6$ inputs, one hidden
    layer of $k=4$ units and an output layer of $K=3$ units, where every hidden and output unit
    carries a bias. Show the count separately for the two weight layers.
  \item \pp{2} A convolutional layer receives a $32\times 32$ image and applies filters of size
    $F=5$ with padding $P=0$ and stride $S=1$. Compute the spatial size of the output feature
    map from $\bigl(W-F+2P\bigr)/S+1$.
\end{enumerate}

\ifdefined\withsolutions
\begin{solutionbox}
\textbf{(a) Forward pass.}
Pre-activations:
\[
  z_1=-2+3\cdot 2-1\cdot 1=-2+6-1=3,\qquad
  z_2=1-2\cdot 2+1\cdot 1=1-4+1=-2 .
\]
ReLU activations:
\[
  A_1=g(3)=\max\{0,3\}=3,\qquad A_2=g(-2)=\max\{0,-2\}=0 .
\]
Output:
\[
  f(x)=-1+2\cdot 3+3\cdot 0=-1+6+0=\mathbf{5}.
\]
The ReLU \emph{switched the second unit off}: its negative pre-activation is clipped to $0$, so
unit~2 contributes nothing for this particular input (though it could for others). This
thresholding is exactly the source of non-linearity: with a purely linear $g$, a network of any
depth would collapse to a single linear function of $x$.

\medskip
\textbf{(b) Softmax and cross-entropy.}
Exponentials $e^{2.0}=7.389$, $e^{1.0}=2.718$, $e^{-1.0}=0.368$ sum to
$7.389+2.718+0.368=10.475$. Softmax probabilities:
\[
  \hat{p}_1=\frac{7.389}{10.475}=0.705,\qquad
  \hat{p}_2=\frac{2.718}{10.475}=0.259,\qquad
  \hat{p}_3=\frac{0.368}{10.475}=0.035,
\]
which sum to $0.999\approx 1$. Cross-entropy loss for the true class~1:
\[
  -\ln\hat{p}_1=-\ln\frac{e^{2}}{10.475}=-\bigl(2-\ln 10.475\bigr)=2.349-2=\mathbf{0.349}.
\]
(Equivalently $-\ln 0.705\approx 0.349$; a perfect prediction $\hat{p}_1=1$ would give loss 0.)

\medskip
\textbf{(c) Parameter count for $6\to 4\to 3$.}
Input $\to$ hidden: $6\cdot 4=24$ weights $+\,4$ biases $=28$.
Hidden $\to$ output: $4\cdot 3=12$ weights $+\,3$ biases $=15$.
Total: $28+15=\mathbf{43}$ parameters.

\medskip
\textbf{(d) Convolution output size.}
\[
  \frac{W-F+2P}{S}+1=\frac{32-5+2\cdot 0}{1}+1=\frac{27}{1}+1=28,
\]
so the output feature map is $\mathbf{28\times 28}$. With no padding a $5\times 5$ filter at
stride $1$ shrinks each spatial dimension by $F-1=4$.
\end{solutionbox}
\fi

% ============================================================
\problem{6: Multiple testing}{20}
{\small\itshape Lecture slides: Chapter 13 --- Multiple Testing (Lecture 12).}\par\vspace{1mm}

\begin{enumerate}[label=(\alph*),leftmargin=7mm,itemsep=2.5mm]
  \item \pp{4} A lab runs $m=10$ \emph{independent} hypothesis tests, each at level
    $\alpha=0.05$, and suppose all $10$ null hypotheses are true. Define the family-wise error
    rate (FWER) and compute it here \emph{(use $(0.95)^{10}=0.5987$)}. Interpret the result in
    one sentence.
  \item \pp{7} A team screens $m=6$ candidate biomarkers; the ordered $p$-values are
    \begin{center}
    \begin{tabular}{lcccccc}
    \toprule
    $j$ & 1 & 2 & 3 & 4 & 5 & 6\\
    \midrule
    $p_{(j)}$ & 0.002 & 0.009 & 0.012 & 0.015 & 0.040 & 0.300\\
    \bottomrule
    \end{tabular}
    \end{center}
    Control the FWER at $\alpha=0.05$ (i) with \emph{Bonferroni} and (ii) with the
    \emph{Holm step-down} procedure. Give the relevant thresholds, state exactly which
    hypotheses each method rejects, and explain why Holm can never reject fewer hypotheses than
    Bonferroni.
  \item \pp{6} Apply the \emph{Benjamini--Hochberg} (BH) procedure at FDR level $q=0.05$ to the
    same six ordered $p$-values ($m=6$; $p_{(1)},\dots,p_{(6)}=0.002,\ 0.009,\ 0.012,\ 0.015,\
    0.040,\ 0.300$): give the thresholds $jq/m$, determine which hypotheses are rejected, and
    state precisely what ``FDR control at $q=0.05$'' guarantees for the resulting set of
    rejections.
  \item \pp{3} A pharmaceutical company runs a \emph{confirmatory} Phase-III trial with $m=4$
    pre-specified primary endpoints; a single significant endpoint could support regulatory
    approval of the drug. Should the company control the FWER or the FDR? Justify your
    recommendation.
\end{enumerate}

\ifdefined\withsolutions
\begin{solutionbox}
\textbf{(a) FWER for $m=10$ independent tests.}
$\text{FWER}=\Pr(\text{at least one Type I error})=\Pr(V\ge 1)$, where $V$ counts the true
nulls that are falsely rejected. With independence and all nulls true,
\[
  \text{FWER}=1-\Pr(\text{no false rejection})=1-(1-0.05)^{10}=1-(0.95)^{10}
  =1-0.5987=\mathbf{0.4013}.
\]
Interpretation: although each individual test uses the $5\%$ level, there is about a $40\%$
chance of at least one false discovery across the ten tests --- testing many hypotheses without
adjustment makes spurious findings very likely.

\medskip
\textbf{(b) Bonferroni and Holm at $\alpha=0.05$.}
\emph{Bonferroni:} reject $H_{(j)}$ iff $p_{(j)}\le\alpha/m=0.05/6=0.00833$.
Only $p_{(1)}=0.002\le 0.00833$ \checkmark; already $p_{(2)}=0.009>0.00833$ $\times$.
$\Rightarrow$ Bonferroni rejects $H_{(1)}$ alone: \textbf{1 rejection}. (Note $p_{(2)}=0.009$ is
below the naive $0.05$ cutoff, yet Bonferroni does not reject it.)

\emph{Holm:} compare $p_{(j)}$ with $\alpha/(m-j+1)$ in order and stop at the first failure:
\begin{center}
\begin{tabular}{lcccccc}
\toprule
$j$ & 1 & 2 & 3 & 4 & 5 & 6\\
\midrule
$p_{(j)}$ & 0.002 & 0.009 & 0.012 & 0.015 & 0.040 & 0.300\\
$\alpha/(m-j+1)$ & 0.00833 & 0.0100 & 0.0125 & 0.0167 & 0.0250 & 0.0500\\
reject? & \checkmark & \checkmark & \checkmark & \checkmark & stop & ---\\
\bottomrule
\end{tabular}
\end{center}
$0.002\le 0.00833$, $0.009\le 0.0100$, $0.012\le 0.0125$, $0.015\le 0.0167$, but
$0.040>0.0250$: stop. $\Rightarrow$ Holm rejects $H_{(1)},H_{(2)},H_{(3)},H_{(4)}$:
\textbf{4 rejections}.

\emph{Why Holm $\supseteq$ Bonferroni:} Holm's thresholds $\alpha/(m-j+1)$ are all $\ge\alpha/m$
(with equality only at $j=1$), so any $p$-value that clears the strict Bonferroni bar also
clears Holm's more lenient later hurdles --- Holm is uniformly at least as powerful, while still
controlling the FWER at $\alpha$ with no independence assumption.

\medskip
\textbf{(c) Benjamini--Hochberg at $q=0.05$.}
Thresholds $jq/m=j\cdot 0.05/6$:
\begin{center}
\begin{tabular}{lcccccc}
\toprule
$j$ & 1 & 2 & 3 & 4 & 5 & 6\\
\midrule
$p_{(j)}$ & 0.002 & 0.009 & 0.012 & 0.015 & 0.040 & 0.300\\
$jq/m$ & 0.00833 & 0.0167 & 0.0250 & 0.0333 & 0.0417 & 0.0500\\
$p_{(j)}\le jq/m$? & yes & yes & yes & yes & \textbf{yes} & no\\
\bottomrule
\end{tabular}
\end{center}
$L=$ largest $j$ with $p_{(j)}\le jq/m$: here $j=5$ (since $0.040\le 0.0417$, while
$0.300>0.0500$). BH rejects \emph{all} $H_{(j)}$ with $j\le L$, i.e.\
$H_{(1)},\dots,H_{(5)}$: \textbf{5 rejections}.
Guarantee: $\text{FDR}=E\bigl[V/\max(R,1)\bigr]\le q=0.05$ --- \emph{on average} at most $5\%$
of the rejected hypotheses are false discoveries. It does \emph{not} promise that at most $5\%$
of these particular $5$ rejections are false, nor does it bound the probability of making any
false rejection (that would be FWER control).

\medskip
\textbf{(d) Confirmatory trial recommendation.}
Control the \textbf{FWER}. Here even a \emph{single} false positive is very costly --- it could
lead to approving an ineffective drug --- so the study must guard against any Type I error, not
merely keep their average proportion small. With only $m=4$ pre-specified endpoints, an FWER
procedure (e.g.\ Holm or Bonferroni at $0.05/4=0.0125$) sacrifices little power. FDR control,
appropriate for large exploratory screens where a few false leads are tolerable, is too
permissive for a confirmatory, decision-driving trial.
\end{solutionbox}
\fi

% ============================================================
\problem{7: Cumulative essentials (Chapters 2--6)}{20}
{\small\itshape Lecture slides: Chapters 2--6 (Lectures 1--8).}\par\vspace{1mm}

\begin{enumerate}[label=(\alph*),leftmargin=7mm,itemsep=2.5mm]
  \item \pp{4} A medical screening classifier is evaluated on $200$ patients (label
    \texttt{1}$=$diseased, \texttt{0}$=$healthy). In \texttt{scikit-learn},
    \texttt{confusion\_matrix(y\_true, y\_pred)} returns the array below (rows $=$ true class,
    columns $=$ predicted class; label order $[0,1]$):
    \begin{center}
    \begin{tabular}{lcc}
    \toprule
     & predicted \texttt{0} & predicted \texttt{1}\\
    \midrule
    true \texttt{0} (healthy)  & $120$ & $30$\\
    true \texttt{1} (diseased) & $10$  & $40$\\
    \bottomrule
    \end{tabular}
    \end{center}
    Read off the counts and compute the overall error rate, the sensitivity, the specificity and
    the precision (positive predictive value) for the positive class \texttt{1}. Which single
    number would trace out the ROC curve if it were varied, and in which direction would
    sensitivity and specificity move?
  \item \pp{4} Consider $K$-nearest-neighbours. (i) Using the training set
    $(x,y)$: $(1,3)$, $(4,6)$, $(6,8)$, $(7,7)$, $(10,12)$, predict the response at $x_0=5$ with $k=3$
    (KNN regression, unweighted average). (ii) Explain how $k$ governs the bias--variance
    trade-off, which extreme ($k=1$ or $k$ large) is the most flexible, and why the
    \emph{training} error of $1$-NN is always $0$.
  \item \pp{4} Explain in one or two sentences what $k$-fold cross-validation does, then give
    one reason to prefer $k=5$ or $10$ over \emph{LOOCV} and one reason to prefer it over a
    \emph{single} validation split.
  \item \pp{2} You have $p=2{,}000$ predictors but expect only a handful to be truly related to
    the response. Ridge or lasso --- which penalty is more appropriate, and where would
    principal-components regression (PCR) sit relative to your choice?
  \item \pp{6} Match each scenario to exactly one method from
    \{\emph{linear regression}, \emph{lasso}, \emph{logistic regression},
    \emph{random forest}\} (each method used exactly once) and justify each choice in one
    sentence:
    \begin{enumerate}[label=(\roman*),nosep,topsep=1mm]
      \item Estimate how much one additional year of schooling adds to annual earnings, with a
        confidence interval, from $n=5{,}000$ survey responses.
      \item Predict a continuous blood-pressure outcome from $p=8{,}000$ genetic markers for
        $n=150$ patients, and report a short list of the relevant markers.
      \item Classify each incoming e-mail as spam or not, and explain to auditors how each
        feature changes the odds of an e-mail being spam.
      \item Predict machine failure from hundreds of mixed sensor variables with strong
        non-linear effects and interactions; only predictive accuracy matters.
    \end{enumerate}
\end{enumerate}

\ifdefined\withsolutions
\begin{solutionbox}
\textbf{(a) Reading the \texttt{confusion\_matrix} output.}
The array is ordered $[0,1]$ with rows $=$ true and columns $=$ predicted, so the diagonal holds
the correct predictions: from the table $TN=120$ and $FP=30$ (true \texttt{0}), $FN=10$ and
$TP=40$ (true \texttt{1}); total $200$, with $50$ diseased and $150$ healthy.
\[
  \text{error rate}=\frac{FP+FN}{200}=\frac{30+10}{200}=\mathbf{0.200},\qquad
  \text{sensitivity}=\frac{TP}{TP+FN}=\frac{40}{50}=\mathbf{0.800},
\]
\[
  \text{specificity}=\frac{TN}{TN+FP}=\frac{120}{150}=\mathbf{0.800},\qquad
  \text{precision}=\frac{TP}{TP+FP}=\frac{40}{70}=\mathbf{0.571}.
\]
Varying the \emph{classification threshold} on the predicted probability traces out the ROC
curve. Lowering the threshold labels more patients ``diseased'', which \emph{raises}
sensitivity (fewer missed cases) but \emph{lowers} specificity (more false alarms); raising it
does the reverse --- the ROC curve captures exactly this trade-off.

\medskip
\textbf{(b) $K$-nearest-neighbours.}
(i) Distances from $x_0=5$ are $|1-5|=4$, $|4-5|=1$, $|6-5|=1$, $|7-5|=2$, $|10-5|=5$; the three
nearest neighbours are $x=4,6,7$ with responses $6,8,7$, so
$\hat{y}(5)=(6+8+7)/3=21/3=\mathbf{7}$.
(ii) Small $k$ gives a very flexible fit: low bias but high variance, since the prediction hangs
on a few nearby points and a wiggly boundary. Large $k$ averages over many neighbours, giving a
smooth, stable fit: high bias but low variance. The most flexible extreme is $\mathbf{k=1}$
(flexibility $\propto 1/k$). The training error of $1$-NN is always $0$ because each training
point's single nearest neighbour is \emph{itself}, so it predicts its own observed response
exactly --- which is why training error is no guide to test performance.

\medskip
\textbf{(c) $k$-fold cross-validation.}
It splits the data into $k$ roughly equal folds; each fold in turn serves as the validation set
while the model is trained on the other $k-1$ folds, and the $k$ resulting errors are averaged.
\emph{vs.\ LOOCV:} the $n$ LOOCV fits use almost identical training sets, so their errors are
highly correlated and their average has \emph{higher variance} than $5$- or $10$-fold CV (and
LOOCV costs $n$ fits, absent shortcut formulas).
\emph{vs.\ a single validation split:} a single split gives a high-variance estimate that
depends heavily on the particular partition and is biased upward because only part of the data
trains the model; $k$-fold CV averages over $k$ folds and trains on $(k-1)/k$ of the data,
giving a more stable, less biased estimate.

\medskip
\textbf{(d) Lasso.}
The $\ell_1$ penalty drives many coefficients \emph{exactly to zero} and thus performs variable
selection --- ideal when the truth is sparse (a few real signals among $2{,}000$ predictors) and
an interpretable small model is wanted. Ridge only shrinks coefficients and keeps all $2{,}000$
predictors, so it tends to lose here. PCR is likewise \emph{not} sparse: it regresses on a few
principal components, so all original variables load on the retained components and no variable
is truly dropped --- like ridge it does not perform selection, so it does not deliver the short
interpretable list that lasso does.

\medskip
\textbf{(e) Matching.}
\begin{itemize}[nosep,leftmargin=5mm]
  \item (i) \textbf{Linear regression}: a quantitative response with $n\gg p$ and the goal of
    \emph{inference} --- an effect size with a confidence interval --- exactly what OLS with
    standard errors provides.
  \item (ii) \textbf{Lasso}: $p\gg n$ makes OLS infeasible; the $\ell_1$ penalty regularises the
    fit and returns a short list of relevant markers.
  \item (iii) \textbf{Logistic regression}: a binary response with interpretable coefficients
    ($e^{\hat\beta_j}$ are odds ratios) --- well suited to explaining odds to auditors.
  \item (iv) \textbf{Random forest}: automatically captures non-linearities and interactions,
    is accurate and robust with many mixed predictors, and interpretability is not required.
\end{itemize}
\end{solutionbox}
\fi

\vspace{6mm}
\begin{center}
\rule{0.35\textwidth}{0.4pt}\\[1mm]
\textit{End of the examination --- good luck!}
\end{center}

\end{document}
